\documentclass[11pt,letterpaper]{amsart}
\usepackage[T1]{fontenc}
\usepackage{lmodern}
\usepackage[margin=1.1in]{geometry}
\usepackage{amsmath,amssymb,amsthm,mathtools}
\usepackage{microtype,booktabs,enumitem,needspace}
\usepackage{xcolor}
\usepackage{hyperref}
\hypersetup{colorlinks=true,linkcolor=blue!45!black,citecolor=blue!45!black,urlcolor=blue!45!black,pdftitle={Hyperrigidity in low dimensions and a dimension-33 counterexample},pdfauthor={Caleb Barnett}}
\setlist[enumerate]{label=\textup{(\roman*)},leftmargin=2em}
\newtheorem{theorem}{Theorem}[section]
\newtheorem{lemma}[theorem]{Lemma}
\newtheorem{proposition}[theorem]{Proposition}
\newtheorem{corollary}[theorem]{Corollary}
\newtheorem{maintheorem}{Theorem}
\renewcommand{\themaintheorem}{\Alph{maintheorem}}
\theoremstyle{definition}
\newtheorem{definition}[theorem]{Definition}
\theoremstyle{remark}
\newtheorem{remark}[theorem]{Remark}
\newtheorem{example}[theorem]{Example}
\numberwithin{equation}{section}
\DeclareMathOperator{\ex}{ex}
\DeclareMathOperator{\dist}{dist}
\DeclareMathOperator{\aff}{aff}
\DeclareMathOperator{\supp}{supp}
\newcommand{\R}{\mathbb R}
\newcommand{\C}{\mathbb C}
\newcommand{\N}{\mathbb N}
\newcommand{\ind}{\mathbf 1}
\newcommand{\B}{\mathcal B}
\title[Hyperrigidity in low dimensions]{Hyperrigidity in low dimensions\\and a dimension-33 counterexample}
\author{Caleb Barnett}
\date{20 October 2026}
\dedicatory{\normalfont\small\textbf{Working draft}\\
Not published or released as a preprint.}
\thanks{Prepared with substantial AI assistance and automated adversarial review. Independent human verification and a complete novelty assessment remain pending; no formal verification is claimed.}
\keywords{Hyperrigidity, unique extension property, compact convex sets, positive operator-valued measures, Riesz energy.}
\begin{document}
\begin{abstract}
We study the least affine dimension of a commutative counterexample to
Arveson's hyperrigidity conjecture. A quadratic estimate for supporting
affine functions relates the unique extension property to inverse-square
energy. Using the normal lift, we prove hyperrigidity for three-dimensional
convex bodies with closed extreme boundary when every extreme point is
exposed or when the boundary is $C^1$. We also modify the Walsh dilation
of Bilich--Scherer to construct a counterexample in affine dimension $33$,
with a $34$-dimensional function system and a diffuse representation.
Finally, we establish two obstructions to smaller versions of that
construction. For the stationary model considered here, arbitrary vector
directions do not yield its
quadratic-distance certificate in horizontal dimension at most four;
more generally, no finite strictly convex graph can repair the stationary
construction in affine dimension four. The latter conclusion requires
no smoothness. The least affine dimension remains between
$3$ and $33$.
\end{abstract}
\maketitle

\section{Introduction}
\label{sec:introduction}

Let $K\subset\R^d$ be a compact convex set with closed extreme boundary
$X=\ex K$. We write $A(K)$ for the complex vector space of continuous
affine functions on $K$ and regard
\[
                         A(K)|_X\subset C(X)
\]
as a function system. A unital representation $\pi:C(X)\to B(H)$ has
the \emph{unique extension property} (UEP) if the only unital completely
positive (UCP) map $\Phi:C(X)\to B(H)$ agreeing with $\pi$ on $A(K)|_X$ is
$\pi$ itself. By Arveson's characterization, the function system is
hyperrigid exactly when every unital representation has this property
\cite{Arv}. We use this formulation throughout.

Each point evaluation at an extreme point has the unique extension
property: a probability measure with that barycenter is a point mass.
The question is whether the same uniqueness survives for diffuse
representations. Bilich and Scherer showed that it need not
\cite{BS}. Their construction raises a geometric question: how small
can the affine dimension of a commutative counterexample be?

The present paper gives positive results in dimension three, an explicit
counterexample in dimension $33$, and structural restrictions on
attempts to bridge the gap. Its first main result is the following.

\Needspace{8\baselineskip}
\begin{maintheorem}\label{intro:positive}
Let $K\subset\R^3$ be a compact convex body with closed extreme
boundary. Then $A(K)|_{\ex K}$ is hyperrigid if either
\begin{enumerate}
\item every extreme point of $K$ is exposed; or
\item $\partial K$ is $C^1$.
\end{enumerate}
In particular, every strictly convex three-dimensional body is covered,
without any smoothness assumption.
\end{maintheorem}

The two hypotheses have different roles. Exposedness supplies an affine
function separating each represented point from every other point.
Smoothness instead supplies a unique normal, which allows a reduction
to planar supporting faces. Neither hypothesis is assumed for general
three-dimensional convex bodies.

The common mechanism is an energy estimate. If $n_x$ and $n_y$ are unit
outward normals at $x,y\in\partial K$, then
\begin{equation}\label{intro:eq:lift}
 0\le\langle n_x,x-y\rangle
 \le\frac14\bigl|(x+n_x)-(y+n_y)\bigr|^2.
\end{equation}
The normal lift $x\mapsto x+n_x$ takes values on
$\partial(K+\overline B)$, where $\overline B$ is the closed unit ball.
In dimension three, every nonzero measure on this parallel surface has
infinite inverse-square energy. On the other hand, a nonzero separated
block of an agreeing positive operator-valued measure would produce a
finite-energy measure. This contradiction proves the required
separation. For a smooth boundary, we first obtain intertwining over
the normal space and only then apply planar hyperrigidity to the
supporting faces. The complete statements and proofs are given in
Theorems~\ref{lb:thm:exposed} and~\ref{lb:thm:c1}.

Our second main result is constructive.

\begin{maintheorem}\label{intro:example}
There is a compact convex body $K\subset\R^{33}$ with closed extreme
boundary $X$, a unital representation $\pi:C(X)\to B(H)$, and a unital
completely positive map $\Phi:C(X)\to B(H)$ such that
\[
                   \Phi|_{A(K)|_X}=\pi|_{A(K)|_X},
                   \qquad \Phi\ne\pi.
\]
The Hilbert space $H$ is separable and infinite dimensional, and the
spectral measure of $\pi$ has no point atoms. Moreover
$\dim_{\C}A(K)=34$.
\end{maintheorem}

The construction is a change of distribution in the Walsh dilation of
Bilich--Scherer, rather than a new dilation mechanism. Theorem~\ref{thm:main}
gives the complete data. The operator identities are proved before the
numerical parameters are chosen; the remaining containment estimate
is a scalar inequality for every infinite digit sequence. A continuous
barycentric kernel then gives the extension, and a quadratic function
provides an explicit positive witness of disagreement.

For comparison, \cite[Corollary 3.3]{BS} proves a construction in
ambient dimension $129$, states an optimization to dimension $57$, and
mentions an unverified dimension-$11$ proposal in Remark~3.4.
Scherer's planar theorem \cite[Theorem 3.8]{planar} excludes dimensions
at most two. Theorem~\ref{intro:example} therefore gives
\begin{equation}\label{eq:dimension-range}
                            3\le d_{\min}\le33,
\end{equation}
where $d_{\min}$ denotes the least affine dimension under the
closed-extreme-boundary hypothesis. The conjectural value four is not
established here. The positive theorem for compact spectrahedra
\cite[Theorem 3.7]{spectra} is a further point of comparison; it applies
to finite self-adjoint affine matrix pencils, not arbitrary projections
of spectrahedra.

The construction obstructions explain why simply changing its visible
geometry does not close this gap. A finite-dimensional signing lemma
rules out the quadratic-distance certificate for every stationary
three-point distribution and every geometric-norm vector encoding in
horizontal dimension at most four, under the nonnegative-potential and
exact-height assumptions of Theorem~\ref{enc:universal}. Thus even arbitrary directions cannot
make that certificate work in affine dimension four or five.
A stronger geometric argument rules out \emph{every finite strictly
convex graph} for the same stationary mechanism in affine dimension
four (Corollary~\ref{prox:cor:walsh}). It uses positive distance levels
of the represented set and requires no differentiability of the graph.
These are obstructions to a specified construction family, not universal
lower bounds on $d_{\min}$.

The general supporting-gap criterion, Theorem~\ref{bar:thm}, is useful
beyond these applications. Appendix~\ref{app:capacity} gives a
source in a four-dimensional convex body whose normal lift supports a
finite-energy measure, but whose representations still have the unique extension
property by a sharper metric. Appendices~\ref{app:heterogeneous}
and~\ref{app:correlated} treat varying and correlated distributions.
Appendix~\ref{app:direct} records two restrictions on direct operator
constructions. Exploratory numerical searches and superseded
certificates are not needed for any proof in this paper.

\section{Representations, measures, and defects}
\label{sec:preliminaries}

All Hilbert spaces are complex. Representations and completely positive
maps are unital unless stated otherwise. A convex body has nonempty
interior in its ambient Euclidean space. We write $\overline B$ for the
closed unit ball and use $|\cdot|$ for Euclidean norms. For a compact
convex set $K$, its support function is
$h_K(u)=\max_{x\in K}\langle u,x\rangle$. An exposed point is the
unique minimizer of an affine function on $K$. A smooth boundary point has a unique unit outward
normal. A strictly convex body has no nondegenerate segment in its
boundary; every boundary point is then exposed.

Whenever $X=\ex K$ is compact, the coordinate functions and the constant
function generate $C(X)$. If $\aff K=\R^d$, their complex linear span
is $A(K)|_X$ and has dimension $d+1$. The dimension of a representation
space is a separate issue.

For a representation $\pi:C(X)\to B(H)$ we write $P$ for its spectral
measure. A UCP map $\Phi$ has a positive operator-valued measure $E$.
One may construct its bounded Borel extension from a Stinespring
representation $\Phi(f)=V^*\rho(f)V$: if $Q$ is the spectral measure of
$\rho$, set $E(B)=V^*Q(B)V$. All operator-valued measures below are
countably additive in the strong operator topology. Saying that a
representation is concentrated on a Borel set $Z$ means $P(Z)=I$;
this is stronger than an assertion about the closure of that set.

\begin{lemma}[Coordinate defects]\label{lem:defect}
Suppose $\Phi=V^*\rho(\,\cdot\,)V$ agrees with $\pi$ on the affine
coordinates $x_1,\ldots,x_d$. Then
\[
 D_j:=\Phi(x_j^2)-\pi(x_j)^2
   =\bigl[(I-VV^*)\rho(x_j)V\bigr]^*
      \bigl[(I-VV^*)\rho(x_j)V\bigr]\ge0.
\]
All $D_j$ vanish if and only if $VH$ reduces $\rho(C(X))$, and this
is equivalent to $\Phi=\pi$.
\end{lemma}
\begin{proof}
Expand the square and use $V^*\rho(x_j)V=\pi(x_j)$.
If it vanishes, $VH$ is invariant for the self-adjoint operator
$\rho(x_j)$ and hence reduces it. The coordinates generate $C(X)$,
so the compression is a representation agreeing with $\pi$ on a
generating set. The converse follows from the same identity.
\end{proof}

\begin{lemma}[Atomic representations]\label{lem:atomic}
Every Hilbert direct sum of point evaluations at points of $X$,
with arbitrary multiplicities, has the unique extension property.
\end{lemma}
\begin{proof}
Write $H=\bigoplus_x H_x$, with $\pi(f)|_{H_x}=f(x)I$.
For a unit vector $h\in H_x$, the scalar measure representing
$f\mapsto\langle\Phi(f)h,h\rangle$ has barycenter $x$ and is therefore
$\delta_x$. Indeed, a probability measure with an extreme barycenter
cannot be split into two parts with different conditional barycenters.
In a Stinespring representation this gives
\[
 \|\rho(f)Vh-f(x)Vh\|^2
       =\langle\Phi(|f-f(x)|^2)h,h\rangle=0.
\]
Thus $\rho(f)V=V\pi(f)$ on each summand and then on their Hilbert direct
sum. This proves equality, including all off-diagonal coefficients.
\end{proof}

In particular, finite-dimensional representations cannot witness a
counterexample. We shall also need a stronger local consequence of
extremality.

\begin{lemma}[Local vanishing of the defect]\label{het:lem:local}
Let $\Phi$ agree with $\pi$ on $A(K)|_X$, and let $D_j$ be the defects
above. If $x$ belongs to the spectral support of $\pi$, then
\[
 \lim_{r\downarrow0}
 \left\|P(B(x,r))\left(\sum_jD_j\right)P(B(x,r))\right\|=0.
\]
Consequently no coordinate defect, nor their sum, is bounded below by
a positive multiple of the identity.
\end{lemma}
\begin{proof}
Any sequence of probability measures on $X$ whose barycenters tend to
$x$ converges weakly to $\delta_x$: every weak limit has the extreme
barycenter $x$. By compactness, their second moments about $x$ tend
to zero uniformly as their barycenters approach $x$.
For a unit vector $h\in P(B(x,r))H$, the scalar measure of $\Phi$
at $h$ has barycenter within $r$ of $x$. Affine agreement gives
\[
 \left\langle\sum_jD_jh,h\right\rangle
 =\int_X|y-x|^2\,d\mu_h(y)
   -\left\langle\pi(|\,\cdot-x|^2)h,h\right\rangle
 \le\int_X|y-x|^2\,d\mu_h(y).
\]
The right side tends to zero uniformly in these unit vectors.
Positivity gives the norm estimate. Every such spectral neighborhood
is nonzero, which proves the last assertion.
\end{proof}

\section{A quadratic affine-barrier criterion}
\label{paper:pos:barriers}

The positive results in this section rest on a common observation.
An affine supporting function can bound the amount of a dilation
that passes from one part of the spectrum to another. If these
bounds vanish quadratically in a suitable metric, a nonzero
off-diagonal effect would produce a scalar measure of finite
inverse-square energy. A critical covering estimate rules out
such a measure. The metric need not be the ambient Euclidean
metric, and this flexibility will be useful below.

We retain the notation of Section~\ref{sec:preliminaries}:
$P$ is the spectral measure of $\pi$, and $E$ is the POVM of
an affine-agreeing UCP map $\Phi$. Their bounded Borel extensions
will be used without changing notation.

\begin{theorem}[Quadratic affine barriers]\label{bar:thm}
Let $\pi:C(X)\to B(H)$ be a unital representation, and let
$Z\subset X$ be Borel with $P(Z)=I$. Suppose that, for each
$x\in Z$, an affine function $b_x$ on $K$ is specified so that
its coefficients depend Borel measurably on $x$ and
\[
 b_x(x)=0,\qquad b_x(y)>0\quad(y\in K\setminus\{x\}).
\]
Suppose also that there are a compact metric space $(T,d_T)$,
a Borel map $\theta:Z\to T$, and a finite constant $M$ such that
\[
 b_x(y)\le M\,d_T(\theta(x),\theta(y))^2
 \qquad(x,y\in Z).
\]
Assume that every nonzero finite positive Borel measure $\mu$
on $Z$ satisfies
\[
 \iint_{Z\times Z}
 d_T(\theta(x),\theta(y))^{-2}\,d\mu(x)d\mu(y)=+\infty,
\]
where the integrand is $+\infty$ when the distance is zero.
Then $\pi$ has the unique extension property relative to
$A(K)|_X$.
\end{theorem}

We first establish the analytic facts needed for the proof.
For a finite positive measure $\nu$ on a metric space, write
\[
 I_2(\nu)=\iint d(u,v)^{-2}\,d\nu(u)d\nu(v).
\]
A Borel set has \emph{zero inverse-square capacity} here if
every nonzero finite positive measure concentrated on it has
infinite $I_2$-energy. The formulation in
Theorem~\ref{bar:thm} is the corresponding statement for the
pushforwards under $\theta$; it does not require a separate
measurability assertion about the image of $\theta$.

\subsection{Energy and localization}

\begin{lemma}[Rank-one domination]\label{lb:lem:rankone}
Let $P$ be a spectral measure on a compact metric space, let
$w$ be a bounded nonnegative Borel function, and suppose
\[
 |\eta\rangle\langle\eta|\le M\int w\,dP.
\]
Then, for $\mu_\eta(D)=\langle P(D)\eta,\eta\rangle$,
\[
                  \int w^{-1}\,d\mu_\eta\le M,
\]
with $w^{-1}=+\infty$ on the zero set of $w$.
\end{lemma}
\begin{proof}
Put $W=\int w\,dP$. For $\epsilon>0$, set
$h=(W+\epsilon I)^{-1}\eta$ and
$t=\langle\eta,h\rangle\ge0$. The assumed domination gives
\[
 t^2\le M\langle Wh,h\rangle
       \le M\langle(W+\epsilon I)h,h\rangle=Mt.
\]
Thus $t\le M$. The spectral theorem and monotone convergence
as $\epsilon\downarrow0$ prove the assertion.
\end{proof}

\begin{lemma}[Critical covering bound]\label{lb:lem:energy}
Let $T$ be a compact metric space. Suppose that, for every
sufficiently small $\epsilon>0$, it can be covered by at most
$M\epsilon^{-2}$ sets of diameter at most $\epsilon$.
Then $I_2(\nu)=+\infty$ for every nonzero finite positive
Borel measure $\nu$ on $T$.
In particular, this conclusion holds for the boundary of
every compact convex body in $\mathbb R^3$.
\end{lemma}
\begin{proof}
The covering sets may be replaced by their closures and then
refined to a Borel partition
$T=\bigsqcup_{j=1}^{N_\epsilon}T_j$, with
$N_\epsilon\le M\epsilon^{-2}$ and
$\operatorname{diam}T_j\le\epsilon$.
If $m=\nu(T)>0$, Cauchy--Schwarz gives
\[
 (\nu\times\nu)\{d(u,v)\le\epsilon\}
 \ge\sum_j\nu(T_j)^2
 \ge\frac{m^2}{N_\epsilon}
 \ge\frac{m^2}{M}\epsilon^2.
\]
The layer-cake formula and Tonelli's theorem therefore imply
\[
 I_2(\nu)
 \ge2\int_0^{\epsilon_0}t^{-3}
       (\nu\times\nu)\{d(u,v)\le t\}\,dt
 \ge\frac{2m^2}{M}\int_0^{\epsilon_0}\frac{dt}{t}
 =+\infty.
\]

For the final assertion, translate the convex body so that
$r\overline B_3\subset K\subset R\overline B_3$ for some
$r,R>0$. Its Minkowski functional $j_K$ is $1/r$-Lipschitz and
is at least $1/R$ on $S^2$. Hence the radial map
$u\mapsto u/j_K(u)$ is Lipschitz from $S^2$ onto $\partial K$.
The sphere admits the required $O(\epsilon^{-2})$ covers,
as follows from finitely many Lipschitz coordinate charts.
The radial map transfers this covering bound to $\partial K$.
\end{proof}

The next lemma isolates the operator argument. Unlike
Theorem~\ref{bar:thm}, it does not require each affine function
to expose a single point.

\begin{lemma}[Local annihilation]\label{paper:pos:localization}
Suppose $\Phi$ and $\pi$ agree on $A(K)|_X$. Let
$F,G\subset X$ be Borel. For each $x\in F$, let $b_x$ be
an affine function nonnegative on $K$. Suppose that
\[
 b_x(y)\ge\delta>0\quad(x\in F,\ y\in G),
 \qquad
 b_x(y)\le M\,d_T(\theta(x),\theta(y))^2
       \quad(x,y\in F),
\]
where $\theta:F\to T$ is Borel and $T$ is compact metric.
If every nonzero finite positive measure on $F$ has infinite
energy in the displayed metric, then $E(G)P(F)=0$.
\end{lemma}
\begin{proof}
Set $A=P(F)E(G)P(F)$. For each fixed $x\in F$, positivity
and affine agreement imply
\[
 0\le A
 \le\delta^{-1}P(F)\Phi(b_x)P(F)
 =\delta^{-1}P(F)\pi(b_x)P(F)
 \le\frac M\delta\int_F
           d_T(\theta(x),\theta(y))^2\,dP(y).
\]
If $A\ne0$, choose a unit vector $h$ such that
$\eta=A^{1/2}h\ne0$. Then
$|\eta\rangle\langle\eta|\le A$ and $\eta\in P(F)H$.
Lemma~\ref{lb:lem:rankone}, applied on $P(F)H$, gives
\[
 \int_Fd_T(\theta(x),\theta(y))^{-2}\,d\mu_\eta(y)
 \le M/\delta\qquad(x\in F).
\]
The same vector $\eta$ works for every $x\in F$. Tonelli's theorem
therefore gives
\[
 \iint_{F\times F}d_T(\theta(x),\theta(y))^{-2}
       \,d\mu_\eta(y)d\mu_\eta(x)
 \le (M/\delta)\|\eta\|^2<\infty,
\]
contrary to the energy hypothesis. Thus $A=0$.
Since $A=(E(G)^{1/2}P(F))^*(E(G)^{1/2}P(F))$, it follows that
$E(G)P(F)=0$.
\end{proof}

We will also use the following elementary passage from compact
sets to the full spectral measure.

\begin{lemma}\label{paper:pos:measure-identification}
Let $P$ be a normalized spectral measure and $E$ a normalized POVM
on a compact metric space. If $E(B)P(A)=0$ whenever $A,B$ are
disjoint compact sets, then $E=P$.
\end{lemma}
\begin{proof}
For an open set $O$, choose compact sets $A_j\uparrow O$.
Applying the hypothesis in both orders to $A_j$ and the
compact set $X\setminus O$, and taking strong limits, gives
\[
 E(X\setminus O)P(O)=0,\qquad
 E(O)P(X\setminus O)=0.
\]
The two effects sum to $I$, so $E(O)P(O)=P(O)$ and
$E(O)P(X\setminus O)=0$. Hence $E(O)=P(O)$.
The scalar matrix-coefficient measures are regular and agree
on open sets; they therefore agree on all Borel sets.
\end{proof}

\begin{proof}[Proof of Theorem~\ref{bar:thm}]
Let $\Phi$ be an agreeing UCP map. Fix disjoint nonempty
compact sets $A,B\subset X$, and put
\[
 F_m=\left\{x\in A\cap Z:
                    \min_{y\in B}b_x(y)\ge1/m\right\}.
\]
The minimum over $B$ is a continuous function of the affine
coefficients, so $F_m$ is Borel. Strict positivity away from
$x$ and compactness of $B$ show that $F_m\uparrow A\cap Z$.
Lemma~\ref{paper:pos:localization} gives $E(B)P(F_m)=0$.
Since $P(Z)=I$, strong convergence yields $E(B)P(A)=0$.
Lemma~\ref{paper:pos:measure-identification} now gives
$E=P$, and hence $\Phi=\pi$.
\end{proof}

\subsection{A graph consequence}

The criterion applies to represented subsets, independently of
the size of the rest of the extreme boundary. The following
consequence is useful when a construction begins with a
convex graph.

\begin{corollary}\label{cap:cor:graph}
Let $B\subset\mathbb R^D$ be compact and convex, and let $g$
be strictly convex and $C^{1,1}$ on an open convex neighborhood
of $B$. Put
\[
 K=\{(a,z):a\in B,\ g(a)\le z\le h\},
 \qquad h>\max_B g,
\]
and assume $\ex K$ is closed. Let $S\subset B$ be compact
and admit covers by $O(\epsilon^{-2})$ sets of diameter at
most $\epsilon$. Every representation concentrated on
\[
                 Z=\{(s,g(s)):s\in S\}
\]
has the unique extension property.
\end{corollary}
\begin{proof}
For $x=(s,g(s))$, take the affine barrier
\[
 b_x(a,z)=z-g(s)-\langle\nabla g(s),a-s\rangle.
\]
Strict convexity makes this nonnegative on $K$, with equality
only at $x$. A Lipschitz bound $L$ for $\nabla g$ on $B$
gives
\[
 b_{(s,g(s))}(t,g(t))\le\tfrac L2|s-t|^2
 \qquad(s,t\in S).
\]
Use $\theta(s,g(s))=s$. Lemma~\ref{lb:lem:energy} and
Theorem~\ref{bar:thm} give the conclusion.
\end{proof}

\section{Normal lifts and exposed points in dimension three}
\label{paper:pos:normal-lifts}

Supporting normals provide a canonical metric for the barrier
criterion. The useful feature is that they can be added to the
boundary points: this places the resulting metric space on the
boundary of a parallel convex body. In dimension three, that
surface has precisely the critical covering bound of
Lemma~\ref{lb:lem:energy}.

For $x\in K$, let
\[
 N_K(x)=\{u\in\mathbb R^d:
               \langle u,y-x\rangle\le0\text{ for every }y\in K\}
\]
be the outward normal cone. For a unit vector $n_x\in N_K(x)$,
put $z_x=x+n_x$. The point $z_x$ lies on
$\partial(K+\overline B_d)$, and its metric projection onto
$K$ is $x$. Indeed, for $y\in K$,
\[
 |x+n_x-y|^2
 =1+|x-y|^2+2\langle n_x,x-y\rangle\ge1,
\]
with equality only when $y=x$. For any two such choices,
\begin{equation}\label{lb:eq:lift}
 0\le\langle n_x,x-y\rangle
 \le\langle n_x-n_y,x-y\rangle
 \le\tfrac14|z_x-z_y|^2.
\end{equation}
The middle inequality uses the supporting property at $y$;
the last follows from $4\langle a,b\rangle\le|a+b|^2$.

\begin{corollary}[Normal-lift criterion]\label{cap:thm}
Let $\pi:C(X)\to B(H)$ have spectral measure concentrated
on a Borel set $Z\subset X$. Suppose a Borel unit normal
$n(x)$ exposes $x$ for every $x\in Z$. If every nonzero finite
positive Borel measure concentrated on
\[
                        L=\{x+n(x):x\in Z\}
\]
has infinite inverse-square energy, then $\pi$ has the unique
extension property.
\end{corollary}
\begin{proof}
The map $x\mapsto x+n(x)$ is Borel and injective. Its image
is Borel as well: if $p_K$ is the continuous metric projection
onto $K$, then
\[
 L=\{z\in\partial(K+\overline B_d):
        p_K(z)\in Z,\ z-p_K(z)=n(p_K(z))\}.
\]
Apply Theorem~\ref{bar:thm} with
$b_x(y)=\langle n(x),x-y\rangle$ and
$\theta(x)=x+n(x)$. The quadratic bound is
\eqref{lb:eq:lift}, and the energy assumption passes to
pushforward measures.
\end{proof}

We need a local version in which the normal need not expose
the source point. This is what allows flat faces in the smooth
case treated in the next section.

\begin{lemma}[Separated normal selections]\label{lb:lem:sep}
Let $K\subset\mathbb R^3$ be a compact convex body, and suppose
$\Phi$ and $\pi$ agree on $A(K)|_X$. Let $F,G\subset X$ be
Borel. If there is a Borel unit normal selection
$x\mapsto n_x$ on $F$ such that
\[
             \langle n_x,x-y\rangle\ge\delta>0
             \qquad(x\in F,\ y\in G),
\]
then $E(G)P(F)=0$.
\end{lemma}
\begin{proof}
Use the affine functions
$b_x(y)=\langle n_x,x-y\rangle$ in
Lemma~\ref{paper:pos:localization}. By \eqref{lb:eq:lift},
their values on $F$ are bounded by one quarter of the squared
distance between the normal lifts. Every nonzero pushforward
measure on these lifts has infinite energy, since the lifts
lie on $\partial(K+\overline B_3)$ and
Lemma~\ref{lb:lem:energy} applies to that surface.
\end{proof}

For exposed points, a separating selection can be made on
successively larger compact sets. We record the elementary
selection argument to make the measurability explicit.

\begin{lemma}[Compact-graph selection]\label{lb:lem:selection}
Let $F,T$ be compact metric spaces, and let
$\Gamma\subset F\times T$ be compact with every fiber
$\Gamma_x$ nonempty. There is a Borel map $s:F\to T$ with
$(x,s(x))\in\Gamma$ for every $x\in F$.
\end{lemma}
\begin{proof}
Fix a countable dense sequence in $T$. At step $j$, choose
the first closed ball of radius $2^{-j}$ centered in that
sequence which meets the current nonempty compact fiber, and
intersect the fiber with that ball. These choices are Borel.
Indeed, for each fixed finite list of earlier choices, the
nonemptiness test is the projection of the original compact
graph intersected with finitely many fixed closed balls;
this projection is compact. The possible earlier lists form
a countable Borel partition, so induction proves the claim.
The remaining fibers are nested nonempty compact sets with
diameters tending to zero. Their unique intersection point
is the limit of the selected centers, and hence depends
Borel measurably on $x$.
\end{proof}

\begin{theorem}[Exposed extreme boundary]\label{lb:thm:exposed}
Let $K\subset\mathbb R^3$ be a compact convex body with
$X=\ex K$ closed. If every point of $X$ is exposed in $K$,
then $A(K)|_X$ is hyperrigid in $C(X)$.
\end{theorem}
\begin{proof}
Let $\pi$ be a unital representation and $\Phi$ an agreeing
UCP map. Fix disjoint nonempty compact sets $A,B\subset X$.
For $m\ge1$, define
\[
 \Gamma_m=\left\{(x,u)\in A\times S^2:
    u\in N_K(x),\quad
    \min_{y\in B}\langle u,x-y\rangle\ge1/m\right\},
 \qquad F_m=\operatorname{pr}_A\Gamma_m.
\]
The graph of the normal cone is closed, and the minimum is
continuous in $(x,u)$. Thus $\Gamma_m$ and $F_m$ are compact.
The sets $F_m$ increase to $A$: an exposing normal at $x$
has a strictly positive minimum on the disjoint compact set
$B$. Lemma~\ref{lb:lem:selection} gives a Borel normal
selection on $F_m$ with gap at least $1/m$ from $B$.
Lemma~\ref{lb:lem:sep} yields $E(B)P(F_m)=0$.
Taking strong limits gives $E(B)P(A)=0$, and
Lemma~\ref{paper:pos:measure-identification} proves
$\Phi=\pi$.
\end{proof}

The selection in this proof depends on the compact target
set $B$. In particular, no continuous choice of exposing
normals is required. For a nonexposed extreme point, the
compact exhaustion can fail when $B$ meets every supporting
face through that point. The next result uses uniqueness of
the normal to deal with such faces instead of separating
their individual points.

\section{Smooth normals and reduction to planar faces}
\label{paper:pos:smooth}

If two smooth boundary points have different normals, the
supporting plane at either point strictly separates it from
the other. The normal-lift estimate therefore identifies the
spectral measure of the normal itself. This is the step that
permits a direct-integral argument: the normal algebra must
first be shown to lie in the multiplicative domain. Agreement on
affine functions alone does not give this. Once the normal is preserved,
the fiber representations are supported on planar exposed faces, where
the planar hyperrigidity theorem applies.

\begin{theorem}[Smooth boundary]\label{lb:thm:c1}
Let $K\subset\mathbb R^3$ be a compact convex body with
$C^1$ boundary and closed extreme boundary $X=\ex K$.
Then $A(K)|_X$ is hyperrigid in $C(X)$.
\end{theorem}

We prove a representation-level statement that also applies
when only the relevant spectral measures are supported on
smooth points. Let
\[
 S=\{x\in X:N_K(x)\cap S^2\text{ consists of one point}\}.
\]
Thus smoothness here refers to uniqueness of the supporting
hyperplane, with no differentiability assumption away from
the represented points.

\begin{proposition}\label{lb:prop:smooth-both}
Let $K\subset\mathbb R^3$ be a compact convex body with
$X=\ex K$ closed. Let $\pi$ and $\rho$ be unital
representations with spectral measures $P$ and $Q$, and
suppose the UCP map $\Phi(f)=V^*\rho(f)V$ agrees with
$\pi$ on $A(K)|_X$. If $P(S)=I$ and $Q(S)=I$, then
$\Phi=\pi$.
\end{proposition}
\begin{proof}
There are two steps. First the energy estimate forces the dilation
isometry to preserve the spectral measure of the normal. This permits
a decomposition into representations on individual supporting faces.
We then apply the planar theorem on those faces.

\emph{The smooth locus.}
The set-valued map $x\mapsto N_K(x)\cap S^2$ has nonempty
compact values and closed graph on $\partial K$. Compactness
of $S^2$ implies that the diameter of these values is upper
semicontinuous. Consequently its singleton locus is Borel,
being the intersection of the relatively open sets where
that diameter is less than $1/j$. On this locus the unique
normal map $n$ is continuous relative to the locus: any
subsequential limit of normals is a normal at the limiting
point, and uniqueness determines that limit. In particular,
$n:S\to S^2$ is Borel.

\emph{Separable reduction.}
It suffices to treat separable representation spaces. If
$\Phi\ne\pi$, choose a witness from a countable dense family
$\{f_j\}\subset C(X)$ and a vector on which the difference
is nonzero. Close its orbit under $\pi(f_j),\Phi(f_j)$ and
their adjoints. The resulting separable subspace $H_0$
reduces both maps and retains the witness. In the given
Stinespring space, the closed span of
$\rho(C(X))VH_0$ is separable, reduces $\rho$, and supplies
a Stinespring representation for the restricted map.
Both spectral concentration assumptions persist under these
restrictions. We may therefore assume that $H$ and the
Stinespring space are separable.

Choose a common scalar control measure $\nu$ for $P$ and $Q$:
all their scalar spectral measures are absolutely continuous with
respect to $\nu$. A positive summable weighted sum of the measures of
countable dense sets of unit vectors gives such a measure. Normalize
it to be a probability measure. It is concentrated on $S$, and
\[
 \nu(B)=0\quad\Longrightarrow\quad P(B)=Q(B)=0.
\]
Only this implication is needed; no inequality of the form
$P(B)\le\nu(B)I$ is asserted.

\emph{Identification of the normal spectral measure.}
For disjoint compact sets $A,B\subset S^2$, put
\[
             S_A=S\cap n^{-1}(A),\qquad
             S_B=S\cap n^{-1}(B).
\]
Inner regularity gives increasing compact sets
$F_j\subset S_A$ and $G_k\subset S_B$ whose unions have
full $\nu$-measure in the respective sets. On
$F_j\times G_k$, the function
$\langle n(x),x-y\rangle$ is continuous and strictly
positive. Indeed, equality would make $n(x)$ a normal at
the smooth point $y$, forcing $n(x)=n(y)$, contrary to the
disjoint normal ranges. Its minimum is therefore positive
whenever both compact sets are nonempty.
Lemma~\ref{lb:lem:sep} gives $E(G_k)P(F_j)=0$.
The sets $S_A\setminus\bigcup_jF_j$ and
$S_B\setminus\bigcup_kG_k$ are $\nu$-null, hence null for both
spectral measures. Consequently
$P(F_j)\to P(S_A)$ and $E(G_k)\to E(S_B)$ strongly. Passing to the
limits one at a time gives
\[
                         E(S_B)P(S_A)=0.
\]

Define measures on $S^2$ by
\[
 P_n(D)=P(S\cap n^{-1}(D)),\qquad
 Q_n(D)=Q(S\cap n^{-1}(D)),\qquad
 E_n(D)=V^*Q_n(D)V.
\]
They are normalized because $P(S)=Q(S)=I$.
The preceding annihilation and
Lemma~\ref{paper:pos:measure-identification} show that
$E_n=P_n$. Compression of a projection to a projection
gives the intertwining relation
\begin{equation}\label{paper:pos:normal-intertwining}
                       Q_n(D)V=VP_n(D)
                       \qquad(D\subset S^2\text{ Borel}).
\end{equation}
Indeed, $V^*Q_n(D)V=P_n(D)$ and $V^*V=I$ give
\[
                    \|Q_n(D)Vh-VP_n(D)h\|^2=0.
\]
Thus the entire normal spectral algebra intertwines through $V$.

\emph{Disintegration.}
The normal map is defined on the possibly nonclosed Borel set $S$,
so we use standard Borel disintegration of measures. Let
$\tau=n_*\nu$, with $\nu$ regarded as a measure on $S$.
The regular conditional probability measures $\nu_\omega$ satisfy
\[
 \nu(B)=\int_{S^2}\nu_\omega(B)\,d\tau(\omega),\qquad
 \nu_\omega(S\cap n^{-1}(\{\omega\}))=1
\]
for $\tau$-almost every $\omega$. These are conditional measures on
the standard Borel space $S$, extended by zero to $X$.
To describe the associated representation decomposition
explicitly, decompose each separable representation into
countably many cyclic representations on
$L^2(X,\mu_j)$, with coordinates acting by multiplication.
Each $\mu_j$ is absolutely continuous with respect to
$\nu$. Choose a Borel density $d_j=d\mu_j/d\nu$.
For each such cyclic component, the identity
\[
 \int_X|f|^2\,d\mu_j
 =\int_{S^2}\!\left(\int_X|f|^2d_j\,d\nu_\omega\right)
       d\tau(\omega)
\]
gives its usual disintegration into the spaces
$L^2(X,d_j\nu_\omega)$. Applying this to $\pi$ and $\rho$
gives a common base measure $\tau$ and decompositions
\[
 H=\int_{S^2}^{\oplus}H_\omega\,d\tau(\omega),\qquad
 L=\int_{S^2}^{\oplus}L_\omega\,d\tau(\omega),
 \qquad
 \pi=\int^{\oplus}\pi_\omega,\quad
 \rho=\int^{\oplus}\rho_\omega.
\]
The countably many conditional measures are finite outside
one null set. On these fields, the normal projections act
as multiplication by $\mathbf 1_D(\omega)$.
Equation~\eqref{paper:pos:normal-intertwining} and the
standard characterization of decomposable operators give
$V=\int^{\oplus}V_\omega\,d\tau(\omega)$.
Uniqueness of decomposable operator fields now turns $V^*V=I$ and
the three coordinate compression identities into
\[
 V_\omega^*V_\omega=I_{H_\omega},\qquad
 V_\omega^*\rho_\omega(x_j)V_\omega=\pi_\omega(x_j)
 \quad(j=1,2,3)
\]
outside one common null set. Zero fibers of $H$ can be ignored.
This is a decomposition of the representations and their isometry;
no pointwise disintegration of an arbitrary UCP map has been assumed.

\emph{Planar fibers.}
For $\omega\in S^2$, let
\[
 F_\omega=\{x\in K:
                 \langle\omega,x\rangle=h_K(\omega)\}.
\]
This exposed face has affine dimension at most two.
The conditional spectral measures above are concentrated
on $S\cap n^{-1}(\{\omega\})\subset F_\omega$.
Consequently, a continuous function vanishing on $X\cap F_\omega$
acts as zero in both fiber representations. They therefore factor
through $C(X\cap F_\omega)$; we keep the same notation for these
factor representations. The face identity
\[
                    X\cap F_\omega=\ex(F_\omega)
\]
shows that this is the function algebra required by the planar
theorem. Indeed, any convex decomposition in $K$ of a point of an
exposed face has both endpoints in that face; extremality in the face
and in $K$ are therefore equivalent there. Its extreme boundary is
compact because
$X$ and $F_\omega$ are closed. There is no need for the
smooth subset of the fiber itself to be closed.

Ambient affine functions restrict onto all affine
functions on $F_\omega$. Hence the planar hyperrigidity
theorem, including its lower-dimensional cases
\cite[Theorem~3.8]{planar}, gives
\[
             V_\omega^*\rho_\omega(f)V_\omega
                  =\pi_\omega(f)
             \qquad(f\in C(X\cap F_\omega))
\]
for almost every $\omega$. Applying this to the
restrictions of a countable dense family in $C(X)$ and
integrating proves $\Phi=\pi$.
\end{proof}

\begin{proof}[Proof of Theorem~\ref{lb:thm:c1}]
On a $C^1$ boundary the supporting unit normal is unique,
so $S=X$. Every representation and every Stinespring
representation are then concentrated on $S$.
Proposition~\ref{lb:prop:smooth-both} applies to every
agreeing UCP map, proving hyperrigidity.
\end{proof}

No positive curvature or strict convexity is needed in
Theorem~\ref{lb:thm:c1}; flat exposed faces are handled by
the planar reduction. The concentration hypothesis on the
Stinespring representation in
Proposition~\ref{lb:prop:smooth-both} is essential to this
proof. A nonsmooth dilation point can have several supporting
normals, so equality of the normal spectral measures does
not follow from the argument above when such points carry
spectral mass.

\section{An explicit counterexample in affine dimension thirty-three}
\label{sec:construction}

\begin{theorem}\label{thm:main}
There is an explicitly specified compact convex body
\(K\subset\mathbb R^{33}\), with \(\aff K=\mathbb R^{33}\) and
compact extreme boundary \(X=\ex K\), for which the affine function
system \(A(K)|_X\) is not hyperrigid in \(C(X)\).
More precisely, there are a unital representation
\(\pi:C(X)\to B(H)\) and a UCP map \(\Psi:C(X)\to B(H)\) such that
\[
 \Psi|_{A(K)|_X}=\pi|_{A(K)|_X},\qquad \Psi(Q)\ne\pi(Q),
 \qquad Q(x,z)=|x|^2.
\]
The Hilbert space \(H\) is separable and infinite dimensional,
the spectral measure of \(\pi\) has no point atoms, and
\(\dim_{\mathbb C}A(K)=34\).
\end{theorem}

We use the Walsh dilation of Bilich--Scherer \cite[Section 3]{BS},
with a different three-point distribution. The operator calculation
is included in full. Its telescoping identity determines the height
of an encoded dilation point. A uniform scalar estimate places all
such points in a truncated epigraph with a strongly convex lower graph. An explicit
boundary kernel then gives the required UCP map.

\subsection{The Walsh compression and its height identity}

Fix \(R>1\) and \(1<b<R^2\), and put
\[
 a=b^{-1},\qquad
 \mu=\frac{b}{2R^2}(\delta_{-R}+\delta_R)
                  +\left(1-\frac b{R^2}\right)\delta_0.
\]
Thus \(\mu\) is symmetric and has second moment \(b\).
In the discrete alphabet \(\{-1,1,-R,0,R,*\}\), with a separate
star symbol, define
\[
 \begin{split}
 G&=\{-1,1\}^{\mathbb N},\qquad Y_0=\{-R,0,R\}^{\mathbb N},\\
 Y_k&=\{-1,1\}^{k-1}\times\{*\}\times
        \{-R,0,R\}^{\{k+1,k+2,\ldots\}}\quad(k\ge1),\\
 Y&=G\cup\bigcup_{k\ge0}Y_k.
 \end{split}
\]
The numerical coordinate \(x_n\) sends the star to zero.
Give \(G\) Haar probability \(m\). Give \(Y_k\) the product
probability \(\mu_k\) with uniform binary prefix, star at \(k\),
and \(\mu\)-distributed tail; on \(Y_0\), use
\(\mu_0=\mu^{\otimes\mathbb N}\).
The space \(Y\) is compact: unbounded star positions can accumulate
only in \(G\), and every fixed branch is compact.
Each coordinate is continuous, and \(Y_0\) is clopen.

Set
\[
 H=L^2(G,m),\qquad
 \mathcal L=\bigoplus_{k\ge0}L^2(Y_k,\mu_k).
\]
For a finite \(F\subset\mathbb N\), write
\(\chi_F=\prod_{n\in F}\varepsilon_n\) and
\(F_{>k}=F\cap\{k+1,k+2,\ldots\}\).
For \(k\notin F\), the functions
\[
                  \chi_F^{(k)}
        =b^{-|F_{>k}|/2}\prod_{n\in F}x_n
\]
form an orthonormal family in the \(k\)-th summand.
Define on the Walsh basis
\begin{equation}\label{eq:V}
 V\chi_F=a^{|F|/2}\chi_F^{(0)}
       +\sqrt{1-a}\sum_{k\in F}
            a^{|F_{>k}|/2}\chi_{F\setminus\{k\}}^{(k)}.
\end{equation}

\begin{lemma}\label{lem:walsh}
The map \(V:H\to\mathcal L\) is an isometry. For bounded Borel
functions on \(Y\), let \(\Phi(f)=V^*M_fV\). Then
\begin{equation}\label{eq:multilinear}
 \Phi\left(\prod_{n\in E}x_n\right)
                  =M_{\prod_{n\in E}\varepsilon_n}
              \qquad(E\subset\mathbb N\text{ finite}).
\end{equation}
Also \(\Phi(\mathbf1_{Y_0})\chi_F=a^{|F|}\chi_F\) and
\(V1=\mathbf1_{Y_0}\).
\end{lemma}
\begin{proof}
On branch \(k\), symmetry and independence make a matrix coefficient
of multiplication by \(\prod_{n\in E}x_n\) vanish unless the two
product labels differ exactly in \(E\) and \(k\notin E\).
In the nonzero case, its value is \(b^{|E_{>k}|/2}\): a tail
coordinate in the multiplier contributes its second moment
divided by \(\sqrt b\), while common coordinates contribute one
after normalization. On branch zero, the same calculation has no
star restriction. Consequently, for \(C_0=F\cap F'\),
\[
 \left\langle M_{\prod_E x_n}V\chi_F,V\chi_{F'}\right\rangle
 =\mathbf1_{\{E=F\triangle F'\}}
 \left(a^{|C_0|}+(1-a)\sum_{k\in C_0}a^{|(C_0)_{>k}|}\right).
\]
The parenthesis is the finite geometric sum
\(a^r+(1-a)\sum_{j=0}^{r-1}a^j=1\).
Taking \(E=\varnothing\) proves isometry on the dense Walsh span;
general \(E\) proves \eqref{eq:multilinear}.
The two final assertions follow by retaining the branch-zero term.
\end{proof}

Let \(\beta\) be an even nonnegative function on the numerical
digits, with \(\beta(\pm1)=0\). Define diagonal contractions
\[
                  A_n\chi_F=a^{|F_{>n}|}\chi_F\qquad(n\ge0).
\]
The parameters \(0<q<1\) and \(c>0\) will be chosen so that
\begin{equation}\label{eq:conditions}
 \begin{aligned}
 \int\bigl(r^2-\beta(r)-1\bigr)\,d\mu(r)
                                      &=c(q^2-1),\\
 a\bigl(R^2-\beta(R)-1\bigr)
       +(1-a)\bigl(-1-\beta(0)\bigr)&=c(q^2-a).
 \end{aligned}
\end{equation}

\begin{lemma}\label{lem:tel}
Under \eqref{eq:conditions},
\begin{equation}\label{eq:tel}
       \Phi\bigl(x_n^2-\beta(x_n)\bigr)-I
                       =c(q^2A_n-A_{n-1})\qquad(n\ge1).
\end{equation}
\end{lemma}
\begin{proof}
Parity makes the left side diagonal. Fix \(F\) and put
\(r=|F_{>n}|\). If \(n\notin F\), the squared weights of the
branches before \(n\) sum to \(a^r\); their expected multiplier
\(x_n^2-\beta(x_n)-1\) is the first left side of
\eqref{eq:conditions}. Later branches contribute zero, since their
\(n\)-th coordinate is binary.
If \(n\in F\), the earlier weights sum to \(a^{r+1}\).
Under the normalized \(x_n^2\)-weighted measure, the coordinate
has magnitude \(R\), so its multiplier is
\(R^2-\beta(R)-1\). The star branch has weight
\((1-a)a^r\) and multiplier \(-1-\beta(0)\).
The two diagonal coefficients are therefore \(a^r\) times the
respective left sides of \eqref{eq:conditions}. These are exactly
the coefficients of the right side of \eqref{eq:tel}.
\end{proof}

Let \(D\) be a positive integer, let
\(j(n)=1+((n-1)\bmod D)\), and define
\begin{equation}\label{eq:encoding}
 \begin{split}
 w_n&=q^{n-1}e_{j(n)},\qquad
 t(y)=\sum_{n\ge1}w_nx_n(y),\qquad s=t|_G,\\
 U(y)&=\sum_{n\ge1}q^{2n-2}\beta(x_n(y)),\\
 z_Y(y)&=|t(y)|^2-U(y)+c\mathbf1_{Y_0}(y).
 \end{split}
\end{equation}
These series converge uniformly and absolutely. All the displayed
functions are continuous, and \(U|_G=0\).
The operator calculation already gives
\begin{equation}\label{eq:height}
                  \Phi(t_j)=M_{s_j},\qquad
                  \Phi(z_Y)=M_{|s|^2}.
\end{equation}
Indeed, all mixed terms in the square expansion of \(|t|^2\)
compress exactly by \eqref{eq:multilinear}. Termwise compression
is valid because \(\sum_n|w_n|<\infty\).
For the diagonal terms, multiplying \eqref{eq:tel} by
\(q^{2n-2}\) gives the norm-convergent telescope
\[
 \sum_{n=1}^Nq^{2n-2}
       \left[\Phi\bigl(x_n^2-\beta(x_n)\bigr)-I\right]
       =c(q^{2N}A_N-A_0)\longrightarrow-cA_0.
\]
This cancels \(c\Phi(\mathbf1_{Y_0})=cA_0\), proving
\eqref{eq:height}. No orthogonality or cancellation of distinct
vectors \(w_n\) is used in this calculation.

\subsection{Numerical choices and uniform containment}

The numerical data are collected below; the other
quantities retain their definitions above.
\begin{center}
\begin{tabular}{c@{\quad}c@{\qquad\qquad}c@{\quad}c}
\toprule
Quantity & Value & Quantity & Value\\ \midrule
\(D\) & \(32\) & \(R\) & \(12\)\\
\(b\) & \(21/2\) & \(q^2\) & \(45026/72405\)\\
\(\beta(0)\) & \(197/200\) & \(\beta(R)\) & \(2417/20\)\\
\(c\) & \(14481/24320\) & Distance coefficient & \(1-10^{-5}\)\\
\bottomrule
\end{tabular}
\end{center}
In particular \(\beta(\pm1)=0\) completes the specification of
\(\beta\). The two scalar identities \eqref{eq:conditions} are
verified exactly: their left sides are respectively
\[
                 -\frac{1441}{6400}=c(q^2-1),\qquad
                   \frac{439}{1400}=c(q^2-a).
\]

For the geometric estimate write \(\rho=q^D\) and
\[
 C_\rho=\left\{\sum_{\ell\ge0}\rho^\ell\varepsilon_\ell:
                              \varepsilon_\ell\in\{-1,1\}\right\}.
\]

\begin{lemma}[Uniform digit estimate]\label{lem:contain}
For every sequence \(d_\ell\in\{-R,-1,0,1,R\}\),
\begin{equation}\label{eq:scalarcontain}
 (1-10^{-5})\dist\left(\sum_{\ell\ge0}\rho^\ell d_\ell,C_\rho\right)^2
             \ge\sum_{\ell\ge0}\rho^{2\ell}\beta(d_\ell).
\end{equation}
\end{lemma}
\begin{proof}
Both sides vanish for an entirely binary sequence. Otherwise let
\(h\) be its first nonbinary position and put
\(\delta=\dist(d_h,\{-1,1\})\in\{1,R-1\}\).
The largest possible difference between a numerical digit and a
binary digit is \(R+1\). Put
\[
 T=\frac{(R+1)\rho}{1-\rho},\qquad
 E=\frac{\beta(R)\rho^2}{1-\rho^2}.
\]
Against any signed expansion, let \(m\le h\) be the first
mismatch. Its leading term dominates the remaining tail and gives
a distance at least
\(\rho^m(|d_m-\varepsilon_m|-T)\).
If \(m=h\), this is at least \(\rho^h(\delta-T)\).
If \(m<h\), it is at least \(\rho^h(R-1-T)\), provided
\(2-T\ge\rho(R-1-T)\). Thus the distance to \(C_\rho\)
is at least \(\rho^h(\delta-T)\).
The right side of \eqref{eq:scalarcontain} is at most
\(\rho^{2h}(\beta(d_h)+E)\), since \(\beta(R)\) is the
largest digit correction.

Here all the required estimates admit rational verification.
The integer inequality
\[
                  1990\cdot45026^{16}<72405^{16}
\]
gives \(\rho<1/1990\), hence
\[
 T<\frac{13}{1989}<1,\qquad
 E<\frac{2417}{20(1990^2-1)}.
\]
Also \(2-13/1989>11/1990\), which implies the earlier-mismatch
comparison \(2-T>\rho(R-1-T)\). Finally, clearing positive
denominators verifies
\[
 \begin{aligned}
 (1-10^{-5})\left(1-\frac{13}{1989}\right)^2
 -\frac{197}{200}-\frac{2417}{20(1990^2-1)}
                                      &>\frac{19}{10000},\\
 (1-10^{-5})\left(11-\frac{13}{1989}\right)^2
 -\frac{2417}{20}-\frac{2417}{20(1990^2-1)}
                                      &>\frac{49}{10000}.
 \end{aligned}
\]
These prove \((1-10^{-5})(\delta-T)^2>\beta(d_h)+E\)
in both cases. The estimates bound the entire infinite tail.
\end{proof}

\subsection{The convex body and its whole extreme boundary}

The represented horizontal set is the Cartesian product
\begin{equation}\label{eq:product}
                 C=s(G)=\prod_{j=1}^Dq^{j-1}C_\rho.
\end{equation}
Define
\begin{equation}\label{eq:g}
 \begin{split}
 g(x)&=|x|^2-(1-10^{-5})\dist(x,C)^2\\
 &=10^{-5}|x|^2+(1-10^{-5})
                  \sup_{u\in C}(2\langle x,u\rangle-|u|^2).
 \end{split}
\end{equation}
The second formula proves strong convexity: it is a positive
quadratic plus a supremum of affine functions.
With
\[
 B=[-R-1,R+1]^D,\qquad M=(R+1)^2D+c+1,
\]
let
\begin{equation}\label{eq:K}
              K=\{(x,z):x\in B,\quad g(x)\le z\le M\}.
\end{equation}
This set is compact and convex. Since
\(g(x)\le|x|^2\le(R+1)^2D<M\) throughout \(B\),
it has nonempty interior in \(\mathbb R^{D+1}\).

\begin{lemma}\label{lem:extreme}
The entire extreme boundary is the compact set
\begin{equation}\label{eq:X}
 X=\{(x,g(x)):x\in B\}
           \cup\{((R+1)\sigma,M):\sigma\in\{-1,1\}^D\}.
\end{equation}
Moreover the continuous maps
\[
 \iota_G(\varepsilon)=(s(\varepsilon),|s(\varepsilon)|^2),
 \qquad \iota_Y(y)=(t(y),z_Y(y))
\]
take their values in \(X\) and \(K\), respectively.
\end{lemma}
\begin{proof}
In a convex decomposition of a lower-graph point, equality in
the convexity inequality forces the horizontal coordinates to
coincide by strict convexity, and then forces the heights to
coincide. Every lower-graph point is therefore extreme.
An intermediate-height point is a nontrivial vertical convex
combination. The top face is a box, whose only extreme points
are its vertices. This proves \eqref{eq:X}, including the absence
of any other extreme points. Compactness follows from continuity
of \(g\) and finiteness of the vertex set.

For the containment assertion,
\( |t_j(y)|\le Rq^{j-1}/(1-\rho)<R+1\).
Also \(U\ge0\), so \(z_Y\le|t|^2+c<M\).
Apply Lemma~\ref{lem:contain} to each residue class of digits
and multiply by \(q^{2(j-1)}\). Squared distance to the product
\eqref{eq:product} is the sum of the squared coordinate distances,
so
\[
                 (1-10^{-5})\dist(t(y),C)^2\ge U(y).
\]
Consequently \(z_Y(y)\ge g(t(y))\), proving \(\iota_Y(Y)\subset K\).
Finally \(g=|\cdot|^2\) on \(C\), so \(\iota_G(G)\subset X\).
\end{proof}

\subsection{The boundary kernel and a positive witness}

For \(w=(x,z)\in K\), put
\[
 r(w)=\frac{z-g(x)}{M-g(x)},\qquad
 p_\sigma(x)=\prod_{j=1}^D\frac{1+\sigma_jx_j/(R+1)}2.
\]
The following probability measure on \(X\) has barycenter \(w\):
\begin{equation}\label{eq:kernel}
 \nu_w=(1-r(w))\delta_{(x,g(x))}
       +r(w)\sum_{\sigma\in\{-1,1\}^D}
                p_\sigma(x)\delta_{((R+1)\sigma,M)}.
\end{equation}
Indeed all weights are nonnegative and sum to one, and
\(\sum_\sigma p_\sigma(x)(R+1)\sigma=x\).
The denominator \(M-g(x)\) is positive on the whole box.
For every \(f\in C(X)\), the function
\(w\mapsto\int f\,d\nu_w\) is continuous.

Define
\begin{equation}\label{eq:maps}
 \pi(f)=M_{f\circ\iota_G},\qquad
 \widehat f(y)=\int_X f\,d\nu_{\iota_Y(y)},\qquad
 \Psi(f)=V^*M_{\widehat f}V.
\end{equation}
The first map is a unital representation. Pointwise integration
against probability measures preserves positivity at every matrix
level; multiplication and isometric compression therefore make
\(\Psi\) UCP. This formula uses the explicit finite kernel
\eqref{eq:kernel}, with no choice of representing measures.

It also gives a representation dilation directly. Index the terms
of \eqref{eq:kernel} by \(J=\{0\}\cup\{-1,1\}^D\), writing
their locations and weights at \(\iota_Y(y)\) as
\(F_\ell(y)\in X\) and \(\omega_\ell(y)\ge0\).
On
\(\widetilde{\mathcal L}
 =\bigoplus_{k\ge0}\bigoplus_{\ell\in J}L^2(Y_k,\mu_k)\),
let \(\rho_X(f)\) multiply component \((k,\ell)\) by
\(f(F_\ell(y))\). The map
\[
              (W\xi)_{k,\ell}(y)
                    =\sqrt{\omega_\ell(y)}\,\xi_k(y)
\]
is an isometry, and
\(\Psi(f)=(WV)^*\rho_X(f)(WV)\).

\begin{proof}[Completion of the proof of Theorem~\ref{thm:main}]
For an affine function, its integral against \(\nu_w\) is its
value at \(w\). Thus \eqref{eq:height} gives
\[
 \Psi(x_j|_X)=M_{s_j}=\pi(x_j|_X),\qquad
 \Psi(z|_X)=M_{|s|^2}=\pi(z|_X).
\]
Together with constants, these coordinates span \(A(K)|_X\).
They are linearly independent because \(K\) has full affine
dimension \(D+1=33\), and restriction to \(X\) is injective
because \(K=\operatorname{conv}X\). Hence the function system
has complex dimension \(D+2=34\).

For \(Q(x,z)=|x|^2\), scalar convexity and the barycenter
property give \(\widehat Q(y)\ge|t(y)|^2\).
The constant Walsh vector has norm one and
\(V1=\mathbf1_{Y_0}\). Independence and zero means remove all
mixed terms from the following scalar integrals:
\[
 \begin{split}
 \langle\Psi(Q)1,1\rangle
       &\ge\int_{Y_0}|t|^2\,d\mu_0=\frac b{1-q^2},\\
 \langle\pi(Q)1,1\rangle
       &=\int_G|s|^2\,dm=\frac1{1-q^2}.
 \end{split}
\]
Consequently
\[
              \langle[\Psi(Q)-\pi(Q)]1,1\rangle
                         \ge\frac{b-1}{1-q^2}>0.
\]
The representation therefore fails the unique extension property.

Finally \(q^D=\rho<1/2\) makes each signed expansion in
\(C_\rho\) injective: the first differing digit exceeds the
entire remaining difference tail. The product encoding \(s\),
and hence \(\iota_G\), is injective. A point inverse image under
\(\iota_G\) is therefore empty or a singleton; Haar measure on
\(G\) assigns it measure zero. This proves that the spectral
measure of \(\pi\) has no point atoms. The countably infinite
Walsh basis proves that \(H\) is separable and infinite dimensional.
\end{proof}

\section{An obstruction for stationary geometric Walsh encodings}
\label{enc:sec}

The construction above separates the operator identities from the
geometric containment problem. We now show that its dimension
cannot be reduced to four by changing only the fixed three-point
distribution, the geometric decay rate, or the directions of the
encoding vectors. The first obstruction concerns the quadratic
distance envelope. The next section extends it, in affine dimension
four, to every finite strictly convex lower graph.

\subsection{A signing estimate}

\begin{lemma}[Signed series in finite dimension]\label{enc:signed}
Let $w_n\in\mathbb R^D$ satisfy $\sum_n|w_n|<\infty$, and let
$L_D$ be the sum of the $D$ largest numbers among
$\{|w_n|^2:n\ge1\}$, padding with zeros if necessary.
There are signs $\varepsilon_n\in\{-1,1\}$ such that
\[
 \left|\sum_{n\ge1}\varepsilon_nw_n\right|^2\le L_D.
\]
\end{lemma}
\begin{proof}
For a finite sequence $w_1,\ldots,w_N$, the polytope
\[
 \left\{a\in[-1,1]^N:\sum_{n=1}^N a_nw_n=0\right\}
\]
is nonempty. At an extreme point $a$, at most $D$ coordinates
lie strictly between $-1$ and $1$: otherwise the corresponding
vectors are linearly dependent and admit a nonzero feasible
perturbation in both directions. Independently round each of
these coordinates to a sign with mean $a_n$, leaving the other
coordinates fixed. Since the centered rounding errors are
independent,
\[
 \mathbb E\left|\sum_{n=1}^N\varepsilon_nw_n\right|^2
 =\sum_{|a_n|<1}(1-a_n^2)|w_n|^2\le L_D.
\]
Some realization therefore has the required bound. Apply this
argument to finite prefixes and extend their signings arbitrarily.
The remaining tail has norm at most $\sum_{n>N}|w_n|\to0$.
Compactness of the sign space and uniform convergence of the
signed series give the assertion for the infinite sequence.
\end{proof}

\subsection{Eliminating all fixed three-point parameters}

In this section, $\gamma$ denotes the squared geometric ratio;
thus $\gamma=q^2$ in the construction above. The numerical tail
law is supported on $\{0,\pm R\}$, with variance $b$ and masses
$1-b/R^2,b/(2R^2),b/(2R^2)$. As before, $G$ is the binary source,
$Y_0$ has only tail coordinates, and $Y_k$ has a binary prefix,
a zero-valued star at $k$, and tail coordinates thereafter.

\begin{theorem}[No stationary geometric distance encoding in
horizontal dimension at most four]\label{enc:universal}
Let
\[
 R>1,\qquad 1<b<R^2,\qquad c>0,\qquad 0<\gamma<1,
 \qquad u,v\ge0.
\]
Suppose the exact height equations hold:
\begin{equation}\label{enc:eq:telescoping}
\begin{aligned}
 \frac b{R^2}(R^2-1-v)
  +\left(1-\frac b{R^2}\right)(-1-u)&=c(\gamma-1),\\
 \frac1b(R^2-1-v)
  +\left(1-\frac1b\right)(-1-u)&=c(\gamma-1/b).
\end{aligned}
\end{equation}
Let $w_n\in\mathbb R^D$ have arbitrary directions and satisfy
$|w_n|^2=\gamma^{n-1}$. Define
\[
 t(y)=\sum_{n\ge1}w_nx_n(y),\qquad C=t(G),\qquad
 U(y)=\sum_{n\ge1}\gamma^{n-1}\beta(x_n(y)),
\]
where $\beta(\pm1)=0$, $\beta(0)=u$, and $\beta(\pm R)=v$.
If $D\le4$, there is no $0<\lambda\le1$ for which
\begin{equation}\label{enc:eq:containment}
 \lambda\dist(t(y),C)^2-U(y)+c\mathbf1_{Y_0}(y)\ge0
 \qquad(y\in Y).
\end{equation}
\end{theorem}
\begin{proof}
Two consequences of Lemma~\ref{enc:signed} suffice. At a star in
the first position followed by zero digits, the encoded point is
zero and $U=u/(1-\gamma)$. Signing the source series bounds its
distance to $C$. On $Y_0$, take digits $R\varepsilon_n$ and compare
with the source point having digits $\varepsilon_n$; the difference
is $(R-1)\sum_n\varepsilon_nw_n$. Thus
\eqref{enc:eq:containment} implies, after discarding $\lambda\le1$,
\begin{equation}\label{enc:eq:necessary}
 u\le1-\gamma^D,\qquad
 v-c(1-\gamma)\le(R-1)^2(1-\gamma^D).
\end{equation}
For $D\le4$ these inequalities remain necessary with $D$ replaced
by four.

Put $x=\gamma^4$, $h=1+u$, and $\Delta=R^2+u-v$.
The height equations give
\begin{equation}\label{enc:eq:elimination}
 c=\frac{\Delta(R^2-b^2)}{R^2(b-1)},\qquad
 c(1-\gamma)=h-\frac b{R^2}\Delta.
\end{equation}
Using the second equality in the second inequality of
\eqref{enc:eq:necessary} yields
\begin{equation}\label{enc:eq:positive-delta}
 \left(1-\frac b{R^2}\right)\Delta
 \ge2(R-1)+(R-1)^2x>0.
\end{equation}
Hence $\Delta>0$, and the first equality of
\eqref{enc:eq:elimination}, together with $c>0$, forces $b<R$.
In particular all subsequent denominators are positive.
Eliminating $c$ also gives
\[
 \frac h\Delta
 =\frac{R^2-b-\gamma(R^2-b^2)}{R^2(b-1)}>0.
\]
Multiply \eqref{enc:eq:positive-delta} by $h/\Delta$ and use
$h\le2-x$, from the first inequality in
\eqref{enc:eq:necessary}. After cancellation,
\begin{equation}\label{enc:eq:reduced}
 (R-1)(2+(R-1)x)
 \left(1-\gamma\frac{R^2-b^2}{R^2-b}\right)
 \le(b-1)(2-x).
\end{equation}

Set $r=R-1>0$, $\alpha=(b-1)/(R-1)\in(0,1)$, and
$A=1-\alpha\in(0,1)$. Clearing the positive denominators in
\eqref{enc:eq:reduced} says that $p(r)\ge0$, where
\[
 p(r)=\gamma(1-\alpha)(2+(1+\alpha)r)(2+rx)
       -(r+2-\alpha)(2(1-\alpha)+(r+\alpha)x).
\]
Its quadratic, linear, and constant coefficients are respectively
\begin{equation}\label{enc:eq:coefficients}
\begin{aligned}
 &x\bigl(\gamma(1-\alpha^2)-1\bigr),\\
 &2\bigl[-\gamma A^2+(2\gamma-1+\gamma x)A-x\bigr],\\
 &-(2-x)A^2+2(2\gamma-1)A-x.
\end{aligned}
\end{equation}
We prove that all three are strictly negative.
The first is immediate. For the second put
$N=2\gamma-1+\gamma^5$. If $N\le0$, negativity is immediate;
if $N>0$, half the coefficient is at most
$N^2/(4\gamma)-\gamma^4$. We claim that
\begin{equation}\label{enc:eq:linear-bound}
 2\gamma-1+\gamma^5<2\gamma^{5/2}\qquad(0<\gamma<1).
\end{equation}
Indeed, setting $t=\sqrt\gamma$, the difference factors as
\[
 t^{10}-2t^5+2t^2-1=(t-1)Q(t),
\]
where
\[
 Q(t)=t^9+t^8+t^7+t^6+t^5-t^4-t^3-t^2+t+1>0
 \qquad(0<t<1).
\]
For $t\le2/3$, this follows from
$t^2+t^3+t^4\le76/81<1$. For $t\ge2/3$, use
\[
 Q(t)=(1-t^2)(1+t)
       +t^4(t+t^2+t^3+t^4+t^5-1),
\]
whose first term is positive and whose second is nonnegative
because $t+t^2\ge10/9$. This proves
\eqref{enc:eq:linear-bound} and negativity of the linear coefficient.

For the constant coefficient in \eqref{enc:eq:coefficients},
$\gamma\le1/2$ is immediate. If $\gamma>1/2$, its maximum over
real $A$ is
\[
 \frac{(2\gamma-1)^2}{2-\gamma^4}-\gamma^4<0.
\]
To verify the strict inequality, factor its numerator as
\[
\begin{split}
 (2\gamma-1)^2-2\gamma^4+\gamma^8
 ={}&(\gamma-1)\bigl(\gamma^7+\gamma^6+\gamma^5+\gamma^4\\
 &\hspace{43mm}-\gamma^3-\gamma^2+3\gamma-1\bigr).
\end{split}
\]
The parenthesis is positive, since
$3\gamma-1-\gamma^2-\gamma^3
=(1-\gamma)(\gamma^2+2\gamma-1)\ge0$ on this interval and the
remaining powers are positive. All coefficients of $p$ are
therefore negative, so $p(r)<0$ for $r>0$, a contradiction.
\end{proof}

Theorem~\ref{enc:universal} permits arbitrary, nonperiodic directions
and every fixed symmetric three-point distribution with nonnegative
potential satisfying the stated height equations. It excludes the distance-envelope realization
in affine dimensions four and five, but is not a universal lower
bound for commutative counterexamples. The remaining possibility
of changing the lower graph requires an operator argument.

\section{Strictly convex graphs in affine dimension four}
\label{prox:sec}

The quadratic envelope is forced by an injective affine defect
when the horizontal dimension is at most three. The proof requires
no smoothness of the lower graph. Away from its quadratic
tangencies, supporting affine functions admit a metric taking
values on a positive distance level. Such levels have the critical
covering bound needed by Theorem~\ref{bar:thm}.

\subsection{Distance levels and supporting affine functions}

\begin{lemma}[Covering a positive distance level]\label{prox:lem:level}
Let $C\subset\mathbb R^D$ be nonempty and compact, and fix
$\varepsilon>0$. The compact set
\[
 S_\varepsilon=\{z:\dist(z,C)=\varepsilon\}
\]
has covers by $O(r^{-(D-1)})$ sets of diameter at most $r$ as
$r\downarrow0$.
\end{lemma}
\begin{proof}
For $z\in S_\varepsilon$, choose a nearest point $c\in C$ and put
$v=(z-c)/\varepsilon$. If $z'\in S_\varepsilon$ and $d=z'-z$,
then $|z'-c|\ge\varepsilon$. Interchanging the two points gives,
for any nearest-point unit vector $v'$ at $z'$,
\begin{equation}\label{prox:eq:normal-bounds}
 v\cdot d\ge-\frac{|d|^2}{2\varepsilon},\qquad
 v'\cdot d\le\frac{|d|^2}{2\varepsilon}.
\end{equation}
Cover $S_\varepsilon$ by finitely many closed spatial sets of
diameter at most $\varepsilon/4$, and the unit sphere by finitely
many closed balls of radius $1/8$ with unit centers $e_j$.
Partition into overlapping patches by requiring membership in
one spatial set and the existence of a nearest-point vector
within $1/8$ of one $e_j$. These patches are compact, since they
are projections of compact sets of nearest-point pairs.

For two points in the same patch, choose witnessing vectors in
\eqref{prox:eq:normal-bounds}. Then
\[
 |e_j\cdot d|\le\frac18|d|+
           \frac{|d|^2}{2\varepsilon}\le\frac14|d|.
\]
Orthogonal projection onto $e_j^\perp$ is therefore injective on
the patch and has a Lipschitz inverse. A grid in its bounded
projected subset of $\mathbb R^{D-1}$ gives the stated covering
bound. Sum over the finitely many patches. For $D=1$, each
patch contains at most one point. No continuity or uniqueness
of nearest-point vectors is needed.
\end{proof}

For a finite convex function $g$ on an open convex set $\Omega$, we
write
\[
 \partial g(s)=\{p:\ g(t)\ge g(s)+\langle p,t-s\rangle
                         \text{ for all }t\in\Omega\}
\]
for its convex subdifferential. The next lemma uses subgradients that
differ from the gradient $2s$ of the quadratic source graph.

\begin{lemma}[A metric for subgradient barriers]\label{prox:lem:metric}
Let $g$ be finite and convex on an open convex set
$\Omega\subset\mathbb R^D$, and let $C\subset\Omega$ be compact
with $g(s)=|s|^2$ on $C$. Suppose $F\subset C$ is Borel and
$p_s\in\partial g(s)$ is a Borel selection satisfying
$0<\delta\le|p_s-2s|\le M<\infty$ on $F$.
Set
\[
 v_s=\frac{p_s-2s}{|p_s-2s|},\qquad
 \varepsilon=\delta/16,\qquad \theta_s=s+\varepsilon v_s.
\]
Then $\theta:F\to S_\varepsilon$ is Borel, with the distance
level taken relative to the whole set $C$, and for a finite
constant $L$ depending only on $\delta,M$,
\begin{equation}\label{prox:eq:barrier-bound}
 0\le|t|^2-|s|^2-p_s\cdot(t-s)
       \le L|\theta_t-\theta_s|^2\qquad(s,t\in F).
\end{equation}
\end{lemma}
\begin{proof}
The subgradient inequality and $g=|\cdot|^2$ on $C$ give
\[
 (p_s-2s)\cdot(t-s)\le|t-s|^2\qquad(s\in F,\ t\in C).
\]
Consequently
\[
 |\theta_s-t|^2\ge\varepsilon^2+
       (1-2\varepsilon/\delta)|t-s|^2.
\]
Thus $s$ is the unique nearest point in $C$ to $\theta_s$, and
$\theta_s\in S_\varepsilon$.

For $s,t\in F$ write $d=t-s$ and $w=v_t-v_s$.
Applying the subgradient inequality in both orders gives
\[
 v_s\cdot d\le|d|^2/\delta,\qquad
 v_t\cdot d\ge-|d|^2/\delta,
\]
so $w\cdot d\ge-2|d|^2/\delta$ and
\begin{equation}\label{prox:eq:metric-control}
 |\theta_t-\theta_s|^2
 =|d+\varepsilon w|^2
 \ge\tfrac34|d|^2+\varepsilon^2|w|^2.
\end{equation}
The barrier in \eqref{prox:eq:barrier-bound} is nonnegative by
convexity. Writing $\alpha_s=|p_s-2s|\le M$, its upper bound is
\[
 |d|^2-\alpha_s v_s\cdot d
 \le(1+M/\delta)|d|^2+M|w|\,|d|.
\]
Young's inequality and \eqref{prox:eq:metric-control} bound this
by $L|\theta_t-\theta_s|^2$ for a finite $L$.
\end{proof}

\subsection{An injective defect forces the quadratic envelope}

\begin{theorem}[An injective defect forces the quadratic envelope]
\label{prox:thm:envelope}
Let $1\le D\le3$, and let $g$ be finite and strictly convex on an
open convex set $\Omega\subset\mathbb R^D$.
Let $K\subset\Omega\times\mathbb R$ be compact and convex, with
$X=\ex K$ compact and $K\subset\{(t,z):z\ge g(t)\}$.
Suppose $C\subset\Omega$ is nonempty and compact, and
\[
 g(s)=|s|^2,\qquad x_s=(s,|s|^2)\in K\qquad(s\in C).
\]
The graph points form a set $E=\{x_s:s\in C\}\subset X$.
Let $\pi:C(X)\to B(H)$ be a unital representation concentrated
on $E$ and having full spectral support there.
Suppose a UCP map $\Psi:C(X)\to B(H)$ agrees with $\pi$ on
affine functions, and for some real affine $h$ the operator
\[
 \mathcal D_h=\Psi(h^2)-\pi(h)^2
\]
is injective. Then $\pi$ is concentrated on points where
$\partial g(s)=\{2s\}$, equivalently where $g$ is differentiable
with $\nabla g(s)=2s$. Moreover,
\begin{equation}\label{prox:eq:envelope}
 2s\in\partial g(s)\quad(s\in C),\qquad
 g(t)\ge|t|^2-\dist(t,C)^2\quad(t\in\Omega).
\end{equation}
\end{theorem}
\begin{proof}
Subdifferentials of a finite convex function on an open convex
domain are nonempty and compact, locally bounded, and have closed
graph. Local boundedness follows by evaluating the supporting
inequality a fixed small distance in every direction inside a
larger compact ball. Hence the subgradient graph over $C$ is
compact, and $|p-2s|\le M$ there for some finite $M$.

For $m\ge1$ let
\[
 \Gamma_m=\{(s,p):s\in C,\ p\in\partial g(s),\
                    |p-2s|\ge1/m\},\qquad
 F_m=\operatorname{pr}_C\Gamma_m.
\]
These sets are compact. The increasing sets $F_m$ exhaust
$\{s\in C:\partial g(s)\ne\{2s\}\}$.
On each nonempty $F_m$, Lemma~\ref{lb:lem:selection} gives a
Borel selector with $1/m\le|p_s-2s|\le M$.
For the graph $E_m=\{x_s:s\in F_m\}$ define
\[
 b_{x_s}(t,z)=z-|s|^2-p_s\cdot(t-s).
\]
The coefficients are Borel, and $b_{x_s}$ is nonnegative on $K$
and vanishes there only at $x_s$. Indeed strict convexity makes
the subgradient inequality strict at every $t\ne s$; at $t=s$,
equality requires $z=g(s)$. This also shows that all graph points
under consideration are exposed extreme points of $K$.

Lemma~\ref{prox:lem:metric}, with $\delta=1/m$, bounds these
barriers on $E_m$ by a constant times squared distance after a
Borel map into $S_{1/(16m)}$. By Lemma~\ref{prox:lem:level}, this
compact level has $O(r^{-(D-1)})$ covers and hence, since $D\le3$,
$O(r^{-2})$ covers. Lemma~\ref{lb:lem:energy} gives infinite
inverse-square energy for every nonzero finite measure pushed
onto it. Theorem~\ref{bar:thm} therefore gives UEP for every
representation concentrated on $E_m$.

Write $P$ for the spectral measure of $\pi$ and $P_m=P(E_m)$.
Restrict $\pi$ and compress $\Psi$ to each nonzero $P_mH$.
These maps still agree on the affine functions, so the preceding
UEP assertion makes them equal. Since $P_m$ commutes with
$\pi(h)$,
\[
 P_m\mathcal D_hP_m=0.
\]
The operator $\mathcal D_h$ is positive by the Schwarz inequality.
Thus $\mathcal D_h^{1/2}P_m=0$, and injectivity forces $P_m=0$.
The Borel set
\[
 E_{\mathrm{bad}}=\bigcup_mE_m
 =\{x_s:s\in C,\ \partial g(s)\ne\{2s\}\}
\]
has spectral projection zero. A finite convex function is
differentiable precisely where its subdifferential is a singleton,
which proves the concentration assertion.

Full spectral support makes the remaining source points dense
in $C$. The closed subgradient graph then gives
$2s\in\partial g(s)$ for every $s\in C$. Finally,
\[
 g(t)\ge\sup_{s\in C}
       \bigl(g(s)+2s\cdot(t-s)\bigr)
       =|t|^2-\dist(t,C)^2,
\]
proving \eqref{prox:eq:envelope}.
\end{proof}

\begin{corollary}[UEP away from quadratic tangencies]
\label{prox:cor:source}
Under the geometric hypotheses of Theorem~\ref{prox:thm:envelope},
every representation concentrated on
$E_{\mathrm{bad}}=\{x_s:s\in C,\ \partial g(s)\ne\{2s\}\}$
has UEP. More generally, for any representation $\pi$ concentrated
on $E$ and any affine-agreeing UCP map $\Psi$,
\[
 \bigl(\Psi(h^2)-\pi(h)^2\bigr)P(E_{\mathrm{bad}})=0
\]
for every real affine $h$.
\end{corollary}
\begin{proof}
Use the increasing compact sets $E_m$ in the preceding proof,
each of which supports only representations with UEP.
If $P(E_{\mathrm{bad}})=I$, then $P(E_m)\uparrow I$ and UEP on
the compressed spaces gives
\[
 P(E_m)\Psi(f)P(E_m)=P(E_m)\pi(f)P(E_m)\qquad(f\in C(X)).
\]
Taking strong limits proves $\Psi=\pi$.
Without that concentration assumption, compressed equality still
gives $P(E_m)\mathcal D_hP(E_m)=0$. Positivity yields
$\mathcal D_hP(E_m)=0$, and the strong limit gives the result.
\end{proof}

\subsection{Application to the Walsh construction}

\begin{lemma}[Injective horizontal defects]\label{reg:lem:injective}
For the stationary three-point Walsh compression, assume
$1<b<R^2$ and write $\Phi(f)=V^*M_fV$.
Each operator $D_n=\Phi(x_n^2)-I$ is positive and injective.
If $t_j=\sum_n(w_n)_jx_n$ is a nonzero linear coordinate with
$\sum_nR|w_n|<\infty$ and $s_j=t_j|_G$, then
\begin{equation}\label{reg:eq:defect-sum}
 \Phi(t_j^2)-M_{s_j}^2=\sum_n(w_n)_j^2D_n
\end{equation}
is also positive and injective. A positive barycentric-kernel
realization with horizontal coordinate $t_j$ has an actual
coordinate defect at least as large as this operator.
\end{lemma}
\begin{proof}
Put $a=1/b$. On the Walsh vector $\chi_F$, let $r=|F_{>n}|$.
If $n\notin F$, the total squared weight of branches before
$n$ is $a^r$, and their tail second moment is $b$; later branches
have binary value squared equal to one. If $n\in F$, the total
weight before $n$ is $a^{r+1}$, their squarefree-product-weighted
second moment is $R^2$, and branch $n$ has weight
$(1-a)a^r$ and value zero. Consequently
\[
 D_n\chi_F=a^r
 \begin{cases}
 b-1,&n\notin F,\\
 R^2/b-1,&n\in F.
 \end{cases}\chi_F.
\]
Every diagonal entry is strictly positive, so $D_n$ is injective;
no uniform positive lower bound is asserted.

Uniform absolute convergence permits expansion of $t_j^2$.
All mixed terms compress exactly by
\eqref{eq:multilinear}, leaving the norm-convergent positive sum
\eqref{reg:eq:defect-sum}. One nonzero coefficient bounds this
sum below by a positive scalar multiple of an injective $D_n$.
Finally, scalar Jensen for a probability measure with barycenter
$(t(y),z(y))$ gives $\int\tau_j^2\,d\nu_y\ge t_j(y)^2$.
Applying the positive compression shows
$\Psi(\tau_j^2)\ge\Phi(t_j^2)$, while affine agreement gives
$\pi(\tau_j)=M_{s_j}$. This proves the last assertion.
\end{proof}

\begin{corollary}[No strictly convex graph in affine dimension four]
\label{prox:cor:walsh}
Take the stationary geometric data of
Theorem~\ref{enc:universal}. Suppose a compact convex body $K$
with compact extreme boundary contains every encoded point
\[
 (t(y),z_Y(y)),\qquad
 z_Y(y)=|t(y)|^2-U(y)+c\mathbf1_{Y_0}(y),\qquad y\in Y,
\]
and lies in the epigraph of a finite strictly convex function
$g$ on an open convex horizontal domain containing $t(Y)$.
Assume $g(s)=|s|^2$ on $C=t(G)$, the signed graph is represented
on $L^2(G,m)$, and a positive barycentric kernel realizes an
affine-agreeing UCP map. Then $D\ge4$.
In particular this mechanism cannot give an affine-dimension-four
counterexample, with or without differentiability of $g$.
\end{corollary}
\begin{proof}
The continuous surjection $t:G\to C$ sends Haar measure to a
measure of full support: the preimage of a nonempty relatively
open set is nonempty and open, and therefore has positive Haar
measure. This does not require injectivity of the encoding.
Some horizontal coordinate has a nonzero coefficient, and
Lemma~\ref{reg:lem:injective} gives an injective defect for the
actual barycentric map.

If $D\le3$, Theorem~\ref{prox:thm:envelope} gives
$g(t)\ge|t|^2-\dist(t,C)^2$. Containment of the encoded points
therefore implies
\[
 \dist(t(y),C)^2-U(y)+c\mathbf1_{Y_0}(y)\ge0\qquad(y\in Y),
\]
contradicting Theorem~\ref{enc:universal} with $\lambda=1$.
\end{proof}

\begin{remark}\label{prox:rem:scope}
Theorem~\ref{prox:thm:envelope} itself does not assume a Walsh
model, stationarity, or geometric vector norms. It uses one
injective affine defect and full spectral support on the
quadratic graph. Its dimension restriction is $D-1\le2$, the
covering exponent of the auxiliary distance level. Accordingly,
the strictly convex graph obstruction applies in affine dimension
four; the stronger affine-dimension-five exclusion in
Theorem~\ref{enc:universal} applies to the distance-envelope
certificate. These are construction-family obstructions, not
universal dimension lower bounds. They leave open constructions
using a different compression, nongeometric vector norms, other
height identities, or geometry outside the strictly convex graph
setting.
\end{remark}

\section{The remaining low-dimensional questions}
\label{sec:questions}

The main results leave both parts of the proposed dimension-four
threshold unresolved. There is no construction here in affine
dimension four, and the positive three-dimensional theorems do not
cover every compact convex body with closed extreme boundary.

In dimension three, one unresolved interaction in the smooth-source
approach is between smooth source points and nonsmooth dilation
points. Different source normals may have supporting faces containing
the same dilation point.
The proof of Proposition~\ref{lb:prop:smooth-both} then gives no
intertwining over the normal space. Distinct supporting faces in
$\R^3$ intersect in at most a line and have at most two common extreme points,
but pointwise affine interpolation on these intersections is not yet
an operator estimate. In particular, a bounded array of interpolation
coefficients need not define a bounded Schur multiplier uniformly in
the number of spectral blocks. If $H_N$ is a symmetric Hadamard matrix
and $C_N$ agrees with it off the diagonal and vanishes on the diagonal,
then
\[
 \|H_N/\sqrt N\|=1,
 \qquad C_N\circ(H_N/\sqrt N)=(J_N-I_N)/\sqrt N,
\]
whose norm is $(N-1)/\sqrt N$. Any argument along this route must
use the actual incidence constraints of the supporting faces, not
only entrywise coefficient bounds.

For dimension four, the signing and proximal arguments rule out the
most direct ways of modifying the stationary geometric Walsh
construction. The expectation obstruction in
Appendix~\ref{app:correlated} rules out the quadratic-distance
containment certificate for every exact height correction for the
first correlated distribution. These results identify failures
of specific mechanisms; they provide no evidence that all
four-dimensional bodies must be hyperrigid. A successful construction
must simultaneously preserve exact compression identities, place the
represented spectrum in the extreme boundary, and retain a nonzero
coordinate defect. None of these requirements follows from the others.

\subsection*{Preparation and verification}
This manuscript was developed and edited with substantial language-model
assistance. Separate automated review tasks rederived the compression
identities, checked the geometric and measure-theoretic arguments,
and tested proposed constructions against the stated hypotheses.
Exact rational arithmetic verifies the numerical certificate in
Section~\ref{sec:construction}. These checks are not independent human
mathematical verification and are not a formal proof certification.
The bibliography identifies the source versions used; the limited
literature comparison does not establish priority for every subsidiary
lemma. The presentation separates complete arguments from the open
questions above.

\appendix
\section{Normal-lift capacity is not a necessary condition}
\label{app:capacity}

The normal lift gives a convenient metric for the supporting-gap
criterion, but it need not give the strongest one. The next example
separates the two tests. Its represented source has the unique extension
property even though its normal lift supports a measure of finite
inverse-square energy.

\begin{theorem}\label{geom:thm}
There are a compact convex body $K\subset\R^4$ with closed extreme
boundary and a compact set $Z\subset\ex K$ such that:
\begin{enumerate}
\item every point of $Z$ is exposed and has a unique unit outward normal;
\item the normal lift of $Z$ supports a probability measure of finite
inverse-square energy;
\item every representation concentrated on $Z$ has UEP relative to $A(K)|_{\ex K}$.
\end{enumerate}
\end{theorem}
\begin{proof}
We construct a thin Cantor set whose normal field records two additional
independent collections of digits. For signs $\omega_n\in\{-1,1\}$,
$n\ge0$, put
\[
 s(\omega)=\left(\sum_{n\ge0}4^{-n}\omega_n,
       10^{-3}\sum_{n\ge0}64^{-n}\omega_n,
       10^{-3}\sum_{n\ge0}128^{-n}\omega_n\right),
 \qquad C=s(\{-1,1\}^{\N_0}).
\]
The first coordinate is an injective Cantor expansion, so $s$ is a
homeomorphism onto $C$. Its first three vector coefficients are
linearly independent by the Vandermonde determinant; hence
$\aff C=\R^3$. Also $C\subset(-2,2)^3$ and $|s|<2$.

If two sequences first differ at index $n$, their first-coordinate
difference is at least $(4/3)4^{-n}$, whereas their second- and
third-coordinate differences are at most
$2\cdot10^{-3}64^{-n}/(1-1/64)$ and
$2\cdot10^{-3}128^{-n}/(1-1/128)$, respectively. In particular,
\begin{equation}\label{app:eq:flat}
 |s_j-t_j|\le\frac{8}{7000}|s_1-t_1|^2
       \quad(j=2,3),\qquad |s-t|^2\le2|s_1-t_1|^2.
\end{equation}
Define continuous functions $a,b:C\to[1,2]$ by
\[
 a(s(\omega))=1+\sum_{k\ge0}2^{-(k+1)}\frac{1+\omega_{3k+1}}2,
 \qquad
 b(s(\omega))=1+\sum_{k\ge0}2^{-(k+1)}\frac{1+\omega_{3k+2}}2.
\]
Set $p(s)=2s+(0,a(s),b(s))$. The preceding estimates imply
\begin{equation}\label{app:eq:interpolation}
 \bigl|\langle p(s)-2s,t-s\rangle\bigr|
       \le\frac{32}{7000}|t-s|^2<\frac12|t-s|^2
       \qquad(s,t\in C).
\end{equation}
Thus the data $g(s)=|s|^2$, $\nabla g(s)=p(s)$ admit the strongly
convex extension
\[
 g(x)=\max_{s\in C}
       \left\{|s|^2+\langle p(s),x-s\rangle+\tfrac12|x-s|^2\right\}.
\]
Indeed, subtracting $|x|^2/2$ leaves a supremum of affine functions.
For $x\in C$, \eqref{app:eq:interpolation} shows that the unique
maximizer is $s=x$, and $g(x)=|x|^2$.

For clarity, this uniqueness also proves differentiability at every
$x\in C$. If $s_h$ maximizes at $x+h$, compactness and continuity
force $s_h\to x$. Let $Q_s$ denote the quadratic inside the maximum.
Then
\[
 Q_x(x+h)-Q_x(x)\le g(x+h)-g(x)
      \le Q_{s_h}(x+h)-Q_{s_h}(x).
\]
Both bounds are $\langle p(x),h\rangle+o(|h|)$.

Choose $M>\max_{[-2,2]^3}g$ and let
\[
 K=\{(x,z):x\in[-2,2]^3,\ g(x)\le z\le M\},
 \qquad Z=\{P_s=(s,|s|^2):s\in C\}.
\]
Strict convexity shows that the full extreme boundary consists of the
lower graph and the eight vertices of the top face, exactly as in
Lemma~\ref{lem:extreme}. It is compact, and $K$ has nonempty interior.
At $P_s$, the supporting plane is strict away from $P_s$, and the unique
outward unit normal is
\[
                         n(s)=\frac{(p(s),-1)}{\sqrt{1+|p(s)|^2}}.
\]

We next show that the lift $P_s+n(s)$ carries finite energy. The paired
lift $(P_s,n(s))$ and the sum lift $P_s+n(s)$ are bi-Lipschitz equivalent:
the supporting-normal inequalities give
\[
 |P_s-P_t|^2+|n(s)-n(t)|^2
 \le |P_s+n(s)-P_t-n(t)|^2
 \le2\bigl(|P_s-P_t|^2+|n(s)-n(t)|^2\bigr).
\]
By \eqref{app:eq:flat}, the paired lift is in turn bi-Lipschitz equivalent
to the coordinates $(s_1,a(s),b(s))$. In one direction all its components
are Lipschitz functions of these three coordinates. In the other,
recover $s$ from the position and then recover $p(s)$ from the normal
by division by its last coordinate, which is bounded away from zero.

Under product measure on the signs, $a$ and $b$ are independent uniform
variables on $[1,2]$. Conditional on the digits defining them, the
remaining contribution to $s_1$ is a translate of
\[
                         \sum_{k\ge0}64^{-k}\omega_{3k}.
\]
Its probability in an interval of length $2r$ is bounded by a fixed
multiple of $r^{1/6}$. To verify this, choose $N$ with
$64^{-N}\le r<64^{-(N-1)}$. An interval of that length meets at most
two level-$N$ cylinders, each of mass $2^{-N}$; the separation follows
by comparing the first differing digit with the geometric tail.
Conditioning and integrating over $a,b$ therefore give a ball bound
of the form $Cr^{2+1/6}$ for $(s_1,a,b)$, uniformly in the center.
The same bound holds on the normal lift by bi-Lipschitz equivalence.
For its pushed-forward probability measure $\nu$, it follows that
\[
 \iint |u-v|^{-2}\,d\nu(u)d\nu(v)
 \le1+2C\int_0^1r^{-3}r^{13/6}\,dr<\infty.
\]

Finally, consider the actual supporting affine functions
\[
 b_{P_s}(x,z)=z-|s|^2-\langle p(s),x-s\rangle.
\]
They expose $P_s$, since
$b_{P_s}(x,z)\ge z-g(x)+|x-s|^2/2$ on $K$. On the represented source,
\eqref{app:eq:flat}--\eqref{app:eq:interpolation} give
\[
                  0\le b_{P_s}(P_t)\le3|s_1-t_1|^2.
\]
The first coordinate lies in the compact interval $[-4/3,4/3]$,
which supports no nonzero measure of finite inverse-square energy.
Theorem~\ref{bar:thm} proves the unique extension property on $Z$.
\end{proof}

\section{Varying distributions and uniformly positive defects}
\label{app:heterogeneous}

The local defect estimate in Lemma~\ref{het:lem:local} imposes a
restriction independent of the choice of convex graph. Here is its
consequence for varying three-point distributions.

At coordinate $n$, let the symmetric distribution on $\{0,\pm R_n\}$
have second moment $b_n$, where $1<b_n<R_n^2$. Put $a_n=b_n^{-1}$.
Use the branch spaces of the Walsh construction, with these varying
tail distributions. Replace the formula for $V$ by
\[
 \begin{split}
 V\chi_F={}&\left(\prod_{n\in F}\sqrt{a_n}\right)\chi_F^{(0)}\\
 &+\sum_{k\in F}\sqrt{1-a_k}
     \left(\prod_{n\in F_{>k}}\sqrt{a_n}\right)
       \chi_{F\setminus\{k\}}^{(k)},
 \end{split}
\]
where the normalized branch monomial is
$\chi_F^{(k)}=(\prod_{n\in F_{>k}}b_n^{-1/2})\prod_{n\in F}x_n$.
The proof of Lemma~\ref{lem:walsh} remains valid because
\[
 \prod_{n\in F}a_n+
 \sum_{k\in F}(1-a_k)\prod_{n\in F_{>k}}a_n=1.
\]
In particular all square-free products compress exactly. For the
coordinate-square defect, parity and the same branch calculation give
\begin{equation}\label{het:eq:Dn}
 \bigl(\Phi(x_n^2)-I\bigr)\chi_F
 =\left(\prod_{i\in F_{>n}}a_i\right)
 \begin{cases}
 b_n-1,&n\notin F,\\
 R_n^2/b_n-1,&n\in F
 \end{cases}
 \chi_F.
\end{equation}

\begin{theorem}\label{het:thm:obstruction}
Assume $\sum_n(b_n-1)<\infty$. Let $w_n\in\R^D$ satisfy
$\sum_n R_n|w_n|<\infty$, with at least one $w_n\ne0$, and put
\[
                       t=\sum_n w_nx_n,
                       \qquad s=\sum_nw_n\varepsilon_n.
\]
There is no compact convex body with compact extreme boundary carrying
an affine-agreeing pair $(\pi,\Psi)$ for which the first $D$ ambient
coordinate functions $\tau_j$ satisfy $\pi(\tau_j)=M_{s_j}$ and
\[
                         \Psi(\tau_j^2)\ge\Phi(t_j^2)
                         \qquad(1\le j\le D).
\]
In particular, no positive barycentric-kernel construction realizes
these data, regardless of its additional coordinates.
\end{theorem}
\begin{proof}
The product $a_*:=\prod_n a_n$ is positive. Equation~\eqref{het:eq:Dn}
gives
\[
 \Phi(x_n^2)-I\ge
 a_*\min\{b_n-1,R_n^2/b_n-1\}I>0.
\]
Uniform convergence and exact compression of mixed products imply
\[
 \Phi(t_j^2)-M_{s_j}^2
       =\sum_n(w_n)_j^2\bigl(\Phi(x_n^2)-I\bigr).
\]
For some $j$, a nonzero summand bounds this defect below by a positive
multiple of $I$. The asserted inequality for $\Psi$ contradicts
Lemma~\ref{het:lem:local}. For a barycentric kernel it follows from
scalar Jensen's inequality for a coordinate square and positivity of
compression.
\end{proof}

The same obstruction applies to finite nonlinear coordinates.

\begin{proposition}\label{het:prop:polynomial}
Under the preceding summability hypothesis, every nonconstant real
multiaffine polynomial $p$ in finitely many coordinates satisfies
\[
                    \Phi(p^2)-M_{p(\varepsilon)}^2\ge c_p I
\]
for some $c_p>0$. It therefore cannot be used as an affine coordinate
of a compact-extreme-boundary barycentric construction.
\end{proposition}
\begin{proof}
Exact square-free compression identifies the defect with $T^*T$, where
$T=M_pV-VM_{p(\varepsilon)}$. Restrict this error to the initial branch
$Y_0$. If $p$ involves the first $N$ variables, its tail factor is
bounded below by the positive product $\prod_{n>N}b_n^{-1/2}$.
It suffices to prove injectivity of the finite-prefix error.

Identify functions on the binary cube with multiaffine polynomials.
Let $\mathcal D$ scale a monomial $x_F$ by $\prod_{i\in F}b_i^{-1}$.
If the error vanishes on a vector represented by $h$, then, with
$q=\mathcal Dh$,
\[
                       p(x)q(x)=\mathcal D(p(\varepsilon)h(\varepsilon))(x)
\]
on the full three-point product grid. Both sides have degree at most
two in each variable, so this is a polynomial identity. The right side
is multiaffine. Whenever $p=p_0+x_ip_1$ with $p_1\ne0$, write
$q=q_0+x_iq_1$; comparison of the $x_i^2$ coefficient gives
$p_1q_1=0$, hence $q_1=0$. Thus $q$ and $h$ are independent of every
active variable of $p$. The identity reduces to
\[
        \bigl[p(x)-p(x_1/b_1,\ldots,x_N/b_N)\bigr]q(x)=0.
\]
The first factor is nonzero, because $p$ is nonconstant and every
$b_i>1$. Hence $q=0$. The finite-prefix error has positive minimum
singular value; tensoring with the bounded-below tail proves the claim.
\end{proof}

\section{Correlated magnitudes and an expectation obstruction}
\label{app:correlated}

Independence of the magnitudes is not essential to the Walsh isometry.
This observation gives a distinct family of exact compression maps.
A simple expectation test, however, excludes an entire range of its
possible geometric realizations.

Fix $b>1$ and $0<\delta<b$. Let $\theta$ be uniform on the circle and
let $\xi_n$ be independent uniform signs, independent of $\theta$. Put
\begin{equation}\label{corr:eq:model}
                   x_n=\xi_n\sqrt{b+\delta\cos(4^n\theta)}.
\end{equation}
These coordinates share their angle and are not independent.
Nevertheless, if $S_n=x_n^2$, then for every finite index set and
integers $0\le k_n\le3$,
\begin{equation}\label{corr:eq:moments}
                   \mathbb E\prod_n S_n^{k_n}
                      =\prod_n\mathbb E S_n^{k_n}.
\end{equation}
Indeed, a nonconstant term in the Fourier expansion has frequency
$\sum_n j_n4^n$ with $|j_n|\le3$. Its largest nonzero term cannot be
cancelled, since $4^m>3\sum_{n<m}4^n$. Thus only products of constant
Fourier coefficients survive. By contrast,
\[
 \mathbb E[(S_n-b)^4(S_{n+1}-b)]=\delta^5/16\ne0,
\]
as follows from the $\cos4t$ coefficient in $\cos^4t$.

On the initial branch use \eqref{corr:eq:model} at every coordinate;
on branch $Y_k$ use the binary prefix, the zero-valued star, and the
same correlated law on its tail. The angle and all prefix and tail
signs are independent. For compactness, use disjoint tagged copies of
the binary symbols and of the compact tail intervals, even where their
numerical values agree. Each fixed branch is a compact image of its
angle and sign space; branches with stars tending to infinity
accumulate only on the binary source. Thus the union is compact,
$Y_0$ is clopen, and every numerical coordinate is continuous.

\begin{proposition}\label{corr:prop:compression}
With these branch laws and the same formula for $V$ as in
\eqref{eq:V}, using $a=b^{-1}$, the map $V$ is an isometry and every
square-free product compresses exactly. Moreover $V1=\ind_{Y_0}$ and
\begin{equation}\label{corr:eq:defect}
 \bigl(\Phi(x_n^2)-I\bigr)\chi_F
  =b^{-|F_{>n}|}
  \begin{cases}
   b-1,&n\notin F,\\
   \delta^2/(2b^2),&n\in F
  \end{cases}\chi_F.
\end{equation}
In particular each coordinate-square defect is positive and injective.
\end{proposition}
\begin{proof}
Independent signs make different monomial labels orthogonal.
Equation~\eqref{corr:eq:moments} gives the same normalized monomial
norms as the independent law. A surviving square-free matrix coefficient
uses only exponents zero and two, so the complete proof of
Lemma~\ref{lem:walsh} applies.
For the square defect, $\mathbb E x_n^4=b^2+\delta^2/2$.
If $n\in F$, the total pre-star branch weight is
$b^{-|F_{>n}|-1}$ and the weighted second moment is
$(b^2+\delta^2/2)/b$; the star contributes zero. Subtracting $I$
gives the second coefficient in \eqref{corr:eq:defect}. If $n\notin F$,
the calculation uses the ordinary second moment $b$, giving the first.
\end{proof}

\begin{theorem}[An obstruction independent of the height correction]
\label{avg:lem:vacuum}
For the correlated model, let $w_n\in\R^D$ be a nonzero absolutely
summable sequence and put
\[
 t=\sum_nw_nx_n,\qquad s=\sum_nw_n\varepsilon_n,
 \qquad C=s(G).
\]
Suppose a continuous height $z$ satisfies
\[
 \Phi(z)=M_{|s|^2},\qquad
 z(y)\ge|t(y)|^2-\dist(t(y),C)^2\quad(y\in Y_0).
\]
Then
\begin{equation}\label{avg:eq:mean}
                   \frac1{2\pi}\int_0^{2\pi}
                     \sqrt{b+\delta\cos\theta}\,d\theta\le1.
\end{equation}
No additive form for $z$, and no geometric decay of $|w_n|$, is assumed.
\end{theorem}
\begin{proof}
Let $L=\sum_n|w_n|^2>0$ and denote the mean in
\eqref{avg:eq:mean} by $m$. Exact height compression and
$V1=\ind_{Y_0}$ give $\mathbb E_{Y_0}z=L$, whereas independent signs
give $\mathbb E_{Y_0}|t|^2=bL$.
For each branch point, compare $t$ with the source point having the
same signs. Averaging first over signs and then over the angle gives
\[
 \mathbb E_{Y_0}\dist(t,C)^2
 \le\mathbb E\left|\sum_nw_n\xi_n
           \bigl(\sqrt{b+\delta\cos(4^n\theta)}-1\bigr)\right|^2
 =(b+1-2m)L.
\]
Containment gives the reverse lower bound $(b-1)L$. Dividing by $L$
proves $m\le1$. Absolute uniform convergence justifies all expansions.
\end{proof}

\begin{corollary}\label{avg:cor:first}
The distribution $b=6/5$, $\delta=1$ admits no such exact height in
any finite horizontal dimension. Suppose a body $K$ satisfies the
geometric hypotheses of Theorem~\ref{prox:thm:envelope}, contains
every encoded point $(t(y),z(y))$, and a positive barycentric kernel
realizes the affine-agreeing map on the signed-source representation.
Then no finite strictly convex graph can repair this model when
$D\le3$.
\end{corollary}
\begin{proof}
The two chords of the square-root function on $[b-\delta,b]$ and
$[b,b+\delta]$ give
\[
 m\ge\sqrt b+
  \frac{\sqrt{b+\delta}+\sqrt{b-\delta}-2\sqrt b}{\pi}
 >\frac{\sqrt b+\sqrt{b+\delta}+\sqrt{b-\delta}}3.
\]
The last inequality uses strict concavity and $\pi>3$. At the stated
parameters, the three square roots exceed $109/100$, $148/100$, and
$44/100$, respectively. Thus $m>301/300>1$.
For the graph assertion, exact mixed products and
\eqref{corr:eq:defect} give an injective defect in any nonzero
horizontal coordinate. A positive barycentric kernel preserves that
lower bound, and the signed source has full spectral support.
Theorem~\ref{prox:thm:envelope} forces the excluded quadratic envelope.
\end{proof}

More generally, the mean-radius condition forces
$1<b<\pi^2/8$ and $\delta^2>4(b-1)$. For the first inequality,
the mean is strictly decreasing in $\delta$ and its limit at
$\delta=b$ is $2\sqrt{2b}/\pi$. For the second, its value is strictly
larger than the endpoint chord average
$(\sqrt{b+\delta}+\sqrt{b-\delta})/2$.
These restrictions leave admissible pairs near $(1,1)$; they do not
exclude the whole correlated family or establish containment for any
remaining pair.

\section{Restrictions on direct operator constructions}
\label{app:direct}

We record two short diagnostics that are independent of the Walsh
construction. They explain why two natural direct constructions cannot
supply the desired diffuse example. Neither is a universal dimension
bound.

\begin{proposition}\label{pert:prop:density}
Suppose an affine-agreeing POVM $E$ satisfies $E(B)\le\nu(B)I$ for
every Borel $B\subset X$, where $\nu$ is a finite positive scalar measure.
Then the spectral support of $\pi$ has at most $\lfloor\nu(X)\rfloor$
points and $E=P$. In particular, an operator-valued density with
integrable operator norm cannot yield a diffuse counterexample.
\end{proposition}
\begin{proof}
For a spectral-support point $x$, choose unit vectors
$h_n\in P(B(x,1/n))H$. Their scalar measures for $E$ are probabilities
dominated by $\nu$, with barycenters tending to $x$. Any weak limit is
$\delta_x$ by extremality, and remains dominated by $\nu$. Therefore
$\nu(\{x\})\ge1$. There are at most $\lfloor\nu(X)\rfloor$ such points,
and Lemma~\ref{lem:atomic} gives equality of the maps.
For a positive density $F$ relative to a scalar measure $\mu$, take
$d\nu=\|F\|\,d\mu$.
\end{proof}

This is an order-domination hypothesis, not merely absolute continuity
with respect to a scalar control measure. It does not automatically
apply to a correction of a spectral measure.

\begin{proposition}\label{nonwalsh:prop:toeplitz}
Let $\rho:C(X)\to L^\infty(\mathbb T)$ be a unital
$*$-representation, and let $P_+$ be the orthogonal projection from
$L^2(\mathbb T)$ onto the scalar Hardy space $H^2$. The compression
\[
                   \Phi(f)=P_+M_{\rho(f)}|_{H^2}
\]
cannot be a nontrivial UCP map agreeing with a representation on
$A(K)|_X$.
\end{proposition}
\begin{proof}
Let $S$ be the unilateral shift on $H^2$ and let $e_0$ be its first
basis vector. For self-adjoint Toeplitz operators $A,B$, put
$u_A=S^*Ae_0$ and $u_B=S^*Be_0$. The identities $S^*AS=A$ and
$SS^*=I-e_0e_0^*$ give
\[
                  S^*[A,B]S-[A,B]=u_Au_B^*-u_Bu_A^*.
\]
If $A,B$ commute and $u_A\ne0$, the two rank-one products agree.
It follows that $u_B=\lambda u_A$ for some real $\lambda$.
For a real symbol $f$, $u_{T_f}$ consists of its strictly positive
Fourier coefficients. The negative coefficients are their conjugates.
Consequently, commuting real scalar Toeplitz symbols are all of the
form $\lambda_jf+c_j$ for one real symbol $f$, unless all are constant.

In an affine-agreeing compression the compressed coordinates commute,
because they equal the coordinates of the representation. Their joint
spectral support therefore lies on an affine line, or is a singleton.
A line meets $\ex K$ in at most two points: every interior point of
its intersection segment with $K$ is nonextreme. The represented
support is thus finite, and Lemma~\ref{lem:atomic} applies.
\end{proof}

In contrast, small norm and low Kraus rank do not themselves impose
rigidity. Any counterexample $\Psi(f)=V^*\rho(f)V$ can be amplified
on $H\oplus L\oplus L\oplus\cdots$ to
$\widehat\pi=\pi\oplus\rho\oplus\rho\oplus\cdots$. The isometry
\[
             W(h,\ell_1,\ell_2,\ldots)=(0,Vh,\ell_1,\ell_2,\ldots)
\]
satisfies
$W^*\widehat\pi(f)W=\Psi(f)\oplus\rho(f)\oplus\rho(f)\oplus\cdots$.
Mixing this compression with the identity gives a nontrivial two-Kraus
channel fixing the affine coordinates and arbitrarily close to the
identity. This is a reformulation of a counterexample, rather than an
independent construction of one.

\begin{thebibliography}{Arv}
\bibitem[Arv]{Arv}
W. Arveson, \emph{The noncommutative Choquet boundary II: Hyperrigidity},
Israel J. Math. \textbf{184} (2011), 349--385.
\url{https://arxiv.org/abs/0810.2751}.

\bibitem[BS]{BS}
B. Bilich and M. Scherer, \emph{Commutative counterexamples to
Arveson's hyperrigidity conjecture}, arXiv:2610.26556v1, 2026.
\url{https://arxiv.org/abs/2610.26556v1}.

\bibitem[Sch1]{planar}
M. Scherer, \emph{The hyperrigidity conjecture for compact convex sets
in $\R^2$}, arXiv:2411.11709v1, 2024.
\url{https://arxiv.org/abs/2411.11709v1}.

\bibitem[Sch2]{spectra}
M. Scherer, \emph{The hyperrigidity conjecture for spectrahedra},
arXiv:2601.16075v1, 2026.
\url{https://arxiv.org/abs/2601.16075v1}.
\end{thebibliography}
\end{document}
